\documentclass[UTF8,a4paper,12pt]{ctexart}
\newcommand{\upcite}[1]{\textsuperscript{\cite{#1}}}

\usepackage{listings}
\usepackage{xcolor}

\definecolor{codegreen}{rgb}{0,0.6,0}
\definecolor{codegray}{rgb}{0.5,0.5,0.5}
\definecolor{codepurple}{rgb}{0.58,0,0.82}
\definecolor{codeblue}{rgb}{0.13,0.13,1}
\definecolor{backcolour}{rgb}{0.96,0.96,0.96}

\lstset{
    language=Python,
    backgroundcolor=\color{backcolour},
    commentstyle=\color{codegreen},
    keywordstyle=\color{codeblue}\bfseries,
    stringstyle=\color{codepurple},
    numberstyle=\tiny\color{codegray},
    basicstyle=\ttfamily\footnotesize,
    breakatwhitespace=false,
    breaklines=true,
    captionpos=b,
    keepspaces=true,
    numbers=left,
    numbersep=5pt,
    showspaces=false,
    showstringspaces=false,
    showtabs=false,
    tabsize=4,
    frame=single,
    rulecolor=\color{black!20}
}

\usepackage[top=2.5cm,bottom=2.5cm,left=2.5cm,right=2.5cm]{geometry}
\usepackage{amsmath,amssymb,bm}
\usepackage{graphicx}
\usepackage{booktabs}
\usepackage{multirow}
\usepackage{array}
\usepackage{float}
\usepackage{longtable}
\usepackage{caption}
\usepackage{subcaption}
\usepackage{enumitem}
\usepackage{hyperref}
\usepackage{xcolor}
\usepackage{setspace}
\usepackage{listings}
\usepackage{fancyvrb}
\usepackage{pdflscape}
\usepackage{adjustbox}

\hypersetup{colorlinks=true,linkcolor=black,urlcolor=blue,citecolor=black}
\captionsetup{font=small,labelfont=bf,labelsep=quad}
\graphicspath{{figures/}{../figures/}}
\setlength{\parindent}{2em}
\setlength{\parskip}{0.35em}
\linespread{1.25}

\lstset{
  basicstyle=\ttfamily\scriptsize,
  breaklines=true,
  frame=single,
  columns=fullflexible,
  keepspaces=true,
  showstringspaces=false,
  literate={μ}{{$\mu$}}1 {³}{{$^3$}}1 {°}{{$^\circ$}}1
}

\newcommand{\fig}[4]{
\begin{figure}[H]
  \centering
  \includegraphics[width=#2\textwidth]{#1}
  \caption{#3}
  \label{#4}
\end{figure}
}

\begin{document}

\begin{titlepage}
\thispagestyle{empty}
\begin{center}
    {\zihao{2}\bfseries 参赛承诺书\par}
    \vspace{2cm}
    \begin{minipage}{0.88\textwidth}
    \zihao{-4}\linespread{1.5}\selectfont
    提交包含此承诺书的PDF文件，表明所有此文件的作者共同承诺：

    \vspace{1em}
    我们完全清楚，在竞赛开始后参赛队员不能以任何方式，包括电话、电子邮件、“贴吧”、QQ群、微信群等，与队外的任何人（包括指导教师）交流、讨论与赛题有关的问题；无论主动参与讨论还是被动接收讨论信息都是严重违反竞赛纪律的行为。

    \vspace{1em}
    我们以中国大学生名誉和诚信郑重承诺，严格遵守竞赛章程和参赛规则，以保证竞赛的公正、公平性。如有违反竞赛章程和参赛规则的行为，我们将受到严肃处理。

    \vspace{1em}
    我们授权北京理工大学数学建模竞赛组织方，可将我们的论文以任何形式进行公开展示（包括进行网上公示，在书籍、期刊和其他媒体进行正式或非正式发表等）。
    \end{minipage}
    \vfill
    {\zihao{-4} 2026年5月\par}
\end{center}
\end{titlepage}

\clearpage
\pagenumbering{arabic}
\setcounter{page}{1}

\begin{center}
{\zihao{3}\bfseries 基于滞后效应与混合机器学习的良乡空气质量联动预测及校园活动建议\par}
\end{center}

\vspace{1.0em}
\noindent{\bfseries 摘要：}
本文针对北京房山良乡地区 2026 年 1 月 1 日至 3 月 31 日逐小时空气质量数据，研究时间、温度、湿度对 PM2.5、PM10、O$_3$、NO$_2$、SO$_2$、CO 六种污染物的影响，并建立面向校园活动安排的空气质量分级预测模型。首先，对原始表格中混入的重复时刻、重复表头和缺失小时进行清洗，得到 2160 条完整小时记录。随后，从描述统计、日内周期、相关矩阵和 0--48 小时滞后相关四个角度分析单一污染物变化规律。结果显示，PM2.5 与温度的最强相关出现在 12 小时滞后处，相关系数为 0.4835；PM10 与温度的最强相关出现在 15 小时滞后处，相关系数为 0.3533；O$_3$ 与温度主要表现为即时正相关，CO 与温度表现为即时负相关。

在联动预测方面，本文构建了时间特征、温湿度当前项及滞后项、污染物自身与交叉滞后项，并比较季节性朴素模型、Ridge 滞后模型、仅含时间温湿度的 VAR（向量自回归）、联合 VAR 以及周期残差混合 VAR。由于数据具有极强的 24 小时周期锁相特性，季节性朴素模型（直接复用前一日同小时值）的平均 MAE 为 0.5544，已具备极高精度。周期残差混合 VAR 的平均 MAE 为 0.7309，虽略高于朴素模型，但其核心价值在于：（1）通过多变量联合建模提取污染物间的短期扰动传递机制，其系数矩阵的主要模式（颗粒物同源性、光化学反应）与大气物理一致，消融实验显示相比独立模型 MAE 改进 71.43\%；（2）保留气象外生变量通道，可响应温湿度非周期性变化；（3）输出联合协方差矩阵，为后续风险评估提供基础。为进一步改进预测性能，本文尝试了稀疏向量自回归模型（Sparse VAR），通过 L1 正则化自动筛选显著的交叉滞后项，相比标准 VAR(2)，平均 MAE 降低 5.65\%（降至 0.7103），平均 $R^2$ 提升至 0.9945。基于稀疏 VAR 的预测结果，依据生态环境部 HJ 633--2026 计算 AQI，测试集 AQI 数值 MAE 为 0.5017。全期 AQI 以轻度污染和良为主，首要污染物主要为 PM2.5。结合第一问发现的 12 小时温度滞后效应，本文为良乡校区提出前瞻性活动建议：当检测到前 12 小时温度持续低于阈值时，预警午后颗粒物累积风险，建议调整户外活动安排。

\vspace{0.8em}
\noindent{\bfseries 关键词：}空气质量；滞后相关；VAR（向量自回归）；稀疏VAR；协整检验；AQI；活动建议

\clearpage

\tableofcontents
\clearpage

\section{问题背景}

空气质量直接影响校园生活与户外活动安排。PM2.5、PM10、O$_3$、NO$_2$、SO$_2$、CO 是空气质量评价中的主要污染物，其浓度既受到排放源变化影响，也受温度、湿度等气象条件制约。温湿度对污染物的影响并不一定只在同一时刻出现，前若干小时的气象状态可能通过扩散、化学反应、颗粒物吸湿增长等过程影响后续浓度。

本题给出了北京房山良乡地区 2026 年 1--3 月逐小时空气质量监测数据，包含六种污染物浓度、温度和相对湿度。本文围绕三个问题展开：第一，解释单一污染物与时间、温度、湿度之间的关系；第二，建立六种污染物协同变化的预测模型；第三，将预测结果转换为空气质量等级，为良乡校区同学安排每天活动提供建议。

\fig{fig01_workflow.pdf}{0.98}{建模流程图}{fig:workflow}

\section{问题分析}

第一问要求说明单一污染物与时间、温度、湿度的关系。由于数据为逐小时序列，单纯计算同步相关不足以揭示滞后效应。因此本文先完成数据清洗，再从月均变化、日内周期、污染物相关矩阵和温湿度滞后相关四个层次给出解释。

第二问要求基于时间、温度、湿度预测六种污染物的协同变化。这里的”协同变化”不解释为污染物之间的强因果关系，而解释为六种污染物在共同时间背景、共同气象条件和相互滞后状态下的联动预测。为确保模型能够输出系数矩阵用于物理机制验证，本文设置可解释的时间滞后特征，并用多模型对比验证预测能力。

第三问要求建立空气质量分级、预测模型。本文依据生态环境部发布的 HJ 633--2026 计算 AQI，并将小时预测结果汇总为上午、下午、晚上三个时段建议。由于校园活动安排更关心可执行建议，本文采用保守原则：每个时段取该时段预测 AQI 的最大值作为活动风险依据。

\section{模型假设}

\begin{enumerate}
  \item 赛题给出的监测数据整体真实可靠，个别重复表头、重复时刻和缺失小时属于记录整理问题，可通过规则清洗修正。
  \item 2026 年 1--3 月良乡地区污染物变化规律在短期内相对稳定，可用于测试期滚动预测。
  \item 温度、湿度是本题给出的主要气象变量；未补充风速、边界层高度、交通流等外部时序，因此相关结论限于题目数据范围。
  \item AQI 分级统一采用 HJ 633--2026 口径，以保证预测和活动建议采用同一现行标准。
  \item 模型计算误差相对于监测数据本身的波动可以忽略。
\end{enumerate}

\section{模型构建与求解}

\subsection{问题一：数据预处理与单污染物规律分析}

\subsubsection{数据清洗}

原始 Excel 文件前 9 行为说明文字，第 10 行为字段名。读取后发现数据区中混入了第二个表头行，且 2026-01-16 00:00 出现重复时刻，2026-02-02 14:00 缺失。本文采用如下处理：

\begin{itemize}
  \item 删除无法解析为时间的记录，包括混入的数据区表头行；
  \item 对重复时刻按同一时间的数值均值合并；
  \item 将时间索引补齐为 1 小时间隔，并对缺失小时做时间线性插值；
  \item 对所有污染物和气象变量转为数值型，检查清洗后缺失值数量。
\end{itemize}

\begin{table}[H]
\centering
\caption{数据清洗报告}
\label{tab:cleaning}
\begin{tabular}{l r}
\toprule
项目 & 数值\\
\midrule
表头后原始行数 & 2162\\
无效时间行数 & 2\\
重复时间戳数 & 1\\
补齐缺失小时数 & 1\\
清洗后小时样本数 & 2160\\
剩余缺失值数 & 0\\
起始时刻 & 2026-01-01 00:00\\
结束时刻 & 2026-03-31 23:00\\
\bottomrule
\end{tabular}
\end{table}

清洗后的数据覆盖 90 天，每天 24 条记录，共 2160 条，与题目说明一致。该处理没有删除正常观测区间，也避免了重复表头和重复时刻影响统计结果。

\subsubsection{描述性统计}

\begin{table}[H]
\centering
\caption{主要变量描述统计}
\label{tab:desc}
\small
\begin{tabular}{lrrrrrrrr}
\toprule
变量 & 最小值 & Q1 & 中位数 & Q3 & 最大值 & 均值 & 标准差 & 偏度\\
\midrule
PM2.5 & 17.0 & 48.0 & 61.0 & 73.0 & 101.0 & 60.51 & 16.01 & -0.11\\
PM10 & 31.0 & 72.0 & 89.0 & 106.0 & 151.0 & 89.47 & 21.82 & 0.06\\
O$_3$ & 5.0 & 31.0 & 41.0 & 61.0 & 87.0 & 43.96 & 19.03 & 0.23\\
NO$_2$ & 13.0 & 41.0 & 53.0 & 67.0 & 76.0 & 52.82 & 15.14 & -0.22\\
SO$_2$ & 2.0 & 2.0 & 3.0 & 3.0 & 4.0 & 2.70 & 0.61 & 0.26\\
CO & 0.3 & 0.5 & 0.8 & 0.9 & 1.2 & 0.74 & 0.28 & 0.04\\
温度 & -7.3 & 2.4 & 7.6 & 11.9 & 20.1 & 7.10 & 6.29 & 0.02\\
湿度 & 39.0 & 65.0 & 84.0 & 91.0 & 100.0 & 78.42 & 14.93 & -0.46\\
\bottomrule
\end{tabular}
\end{table}

由表 \ref{tab:desc} 可见，PM10 的均值和波动幅度最大，PM2.5 次之；SO$_2$ 全期浓度低且波动很小；温度覆盖冬末到初春的升温过程，湿度整体偏高。偏度绝对值均不大，说明数据中没有极端偏斜变量，但污染物之间的波动尺度差异明显，后续建模需要分别评价各污染物预测误差。注：除 CO（mg/m$^3$）外，其余污染物单位为 $\mu$g/m$^3$，温度单位为 $^\circ$C，湿度单位为 \%。

\fig{fig02_pollutant_timeseries.pdf}{0.98}{污染物小时浓度与日均值}{fig:pollutant_series}

图 \ref{fig:pollutant_series} 显示 PM2.5 和 PM10 在 1--3 月整体呈上升趋势并叠加明显的日内与周度周期波动，小时数据呈现密集的"毛刺"状，而日均值曲线则平滑揭示了低频演变趋势。O$_3$ 在 2 月后显著升高，符合春季日照时间变长、光化学反应增强的基本规律；NO$_2$ 在 2 月中旬出现一个明显的凸起峰值，随后在 3 月有所回落，可能与该时段特定的静稳天气过程有关；SO$_2$ 和 CO 的变化幅度极小，日均线几乎呈水平直线，说明本地固定排放源在这三个月内非常稳定。

\fig{fig03_meteorology.pdf}{0.90}{逐日气象条件变化}{fig:meteo}

图 \ref{fig:meteo} 中温度由 1 月低温逐步升至 3 月较高水平，湿度长期维持在较高区间。温度上升可能促进 O$_3$ 生成，也可能改变扩散条件和颗粒物吸湿增长；湿度高时颗粒物二次生成和吸湿增长更容易发生，但其影响可能与温度、日周期共同作用。

\fig{fig04_pollutant_distribution.pdf}{0.82}{污染物浓度分布}{fig:distribution}

图 \ref{fig:distribution} 显示 PM10、PM2.5、O3 和 NO2 的箱体较长且上须延伸较远，离散程度较大，是空气质量波动的主要来源。值得注意的是，PM10 的上方存在少量超出上须的离群点，表明监测期内曾发生过短期的粗颗粒物异常升高事件（如局部扬尘或传输过程）；而 SO$_2$ 箱体极窄且无离群点，进一步印证了其在本数据中对 AQI 等级变化贡献有限。

\subsubsection{时间周期与相关结构}

为形式化刻画污染物浓度随时间的变化结构，本文采用加性分解模型将小时序列 $Y_t$ 拆解为趋势项、日周期项和残差项：
\begin{equation}
\label{eq:decomp}
Y_t = T_t + S_t + R_t,
\end{equation}
其中 $T_t$ 为长期趋势（通过 168 小时即一周的中心移动平均提取），$S_t$ 为日周期成分（各小时在扣除趋势后的均值），$R_t=Y_t-T_t-S_t$ 为残余波动。趋势项反映 1--3 月的季节性演变，周期项反映日内排放与光化学节律，残差项反映天气过程和随机扰动。图 \ref{fig:pollutant_series} 中的日均值曲线大致对应 $T_t$，图 \ref{fig:diurnal} 对应 $S_t$ 的归一化结果。

\fig{fig05_diurnal_heatmap.pdf}{0.86}{污染物平均日内变化热图}{fig:diurnal}

图 \ref{fig:diurnal} 表明六种污染物存在显著且互异的日内结构。O$_3$ 的归一化均值在 14:00--15:00 达到全天最高峰，与午后太阳辐射强度最大高度吻合；NO$_2$ 与 CO 在 07:00--08:00 出现明显的早高峰，对应早间通勤排放累积与边界层尚未完全抬升，随后在午后降至谷底；PM2.5 与 PM10 的日内分布相对平坦，但在夜间 22:00--02:00 略有抬升，体现了夜间边界层压缩导致的颗粒物持续累积效应。

\fig{fig06_correlation_matrix.pdf}{0.72}{相关系数矩阵}{fig:corr}

从图 \ref{fig:corr} 可见，PM2.5 与 PM10 的相关系数为 0.982，说明两类颗粒物具有高度同步性。PM10 与 NO$_2$、SO$_2$、CO 也呈正相关，体现出共同排放或共同扩散条件影响。O$_3$ 与 CO、SO$_2$、PM10 呈明显负相关，说明臭氧与一次污染物在日内变化机制上存在差异。这一负相关性可以从光化学机制解释：白天光照增强时，NO$_2$ 光解生成 O$_3$（$\text{NO}_2+h\nu\to\text{NO}+\text{O}$，$\text{O}+\text{O}_2\to\text{O}_3$），同时 O$_3$ 通过与 NO 反应消耗（$\text{O}_3+\text{NO}\to\text{NO}_2+\text{O}_2$）。在日内尺度上，O$_3$ 峰值对应 NO$_2$ 谷值和 CO 谷值（白天边界层抬升也稀释了一次污染物），因此呈现负相关。这不是因果关系，而是共同受日照和边界层调控的协同表现。

\subsubsection{温湿度滞后效应}

为刻画滞后效应，设污染物浓度为 $Y_t$，温度或湿度变量为 $X_t$，滞后 $k$ 小时的相关系数定义为
\begin{equation}
\label{eq:lag_corr}
r_k=\frac{\sum_t (Y_t-\bar{Y})(X_{t-k}-\bar{X}_k)}
{\sqrt{\sum_t (Y_t-\bar{Y})^2}\sqrt{\sum_t (X_{t-k}-\bar{X}_k)^2}},\quad k=0,1,\ldots,48.
\end{equation}
其中 $r_k$ 衡量前 $k$ 小时温湿度与当前污染物浓度之间的线性关系。
（注：当采用重叠区间计算时 $\bar{X}_k$ 随 $k$ 略有变化，本文实现中取全序列均值 $\bar{X}$ 以简化计算，两者差异在本数据尺度下可忽略。）

\begin{table}[H]
\centering
\caption{温湿度最强滞后相关汇总}
\label{tab:lag}
\small
\begin{tabular}{llrr}
\toprule
污染物 & 气象因子 & 最佳滞后（小时） & 相关系数\\
\midrule
PM2.5 & 温度 & 12 & 0.4835\\
PM2.5 & 湿度 & 0 & 0.1727\\
PM10 & 温度 & 15 & 0.3533\\
PM10 & 湿度 & 11 & 0.1024\\
O$_3$ & 温度 & 0 & 0.4040\\
O$_3$ & 湿度 & 35 & -0.1115\\
NO$_2$ & 温度 & 10 & 0.1840\\
NO$_2$ & 湿度 & 0 & 0.1394\\
SO$_2$ & 温度 & 29 & 0.1678\\
SO$_2$ & 湿度 & 12 & 0.1931\\
CO & 温度 & 0 & -0.5210\\
CO & 湿度 & 1 & 0.2376\\
\bottomrule
\end{tabular}
\end{table}

\fig{fig07_lag_correlation.pdf}{0.98}{气象变量滞后相关热图}{fig:lag}

表 \ref{tab:lag} 与图 \ref{fig:lag} 显示，PM2.5、PM10 对温度的最强相关并非即时出现，而分别出现在 12 小时和 15 小时附近。这说明颗粒物浓度受气象条件影响时存在累积和滞后过程。O$_3$ 与温度的相关性在即时最强，符合臭氧生成对当时光热条件更敏感的特点。CO 与温度即时负相关，可能反映升温时扩散条件改善或燃烧排放减弱。

这种滞后效应可以从大气边界层动力学角度进行解释：通常在夜间至清晨，地表辐射冷却导致逆温层形成，边界层高度被压缩至数百米，污染物在垂直方向无法扩散，持续累积；而午后（约滞后12-15小时）随着太阳辐射增强，边界层打破并大幅抬升，此前累积的高浓度气团被垂直混合稀释。因此，当前时刻观测到的颗粒物浓度，实际上是前半夜低温累积过程与当前扩散条件的综合结果，从而表现出与十余小时前温度的显著滞后相关性。

湿度对污染物的滞后影响机制与温度有所不同。高湿度主要通过两条路径影响颗粒物浓度：一是通过气粒转化过程（如 SO$_2$、NO$_x$ 氧化生成硫酸盐、硝酸盐）促进二次颗粒物生成，该过程通常需要数小时至十余小时，因此 SO$_2$ 与湿度的最强相关出现在 12 小时滞后处；二是通过颗粒物吸湿增长改变粒径分布和光学特性，该过程接近即时，因此 PM2.5 与湿度的最强相关出现在 0 小时（即时）。表 \ref{tab:lag} 中 PM10 与湿度的 11 小时滞后可能反映了粗颗粒物在湿润条件下经历沉降、再悬浮的较慢过程。O$_3$ 与湿度的 35 小时负相关则可能与高湿度条件下光化学反应受抑制有关：水汽通过竞争羟基自由基（$\cdot$OH）和促进云雾沉降等途径削弱臭氧生成，该效应在累积一天半后最为显著。

\fig{fig08_3d_pm25_meteorology.pdf}{0.72}{PM2.5 温湿度响应曲面（探索性分析）}{fig:3d}

图 \ref{fig:3d} 展示了 PM2.5 与温湿度的三维响应关系的探索性分析。\textbf{该曲面由随机森林模型拟合生成}，用于直观展示数据中温湿度与 PM2.5 的联合分布特征。曲面呈现出明显的非线性特征：在湿度较低区间，曲面相对平缓；在湿度较高区间，曲面出现明显抬升。红色散点为实际观测值，密集分布在拟合曲面周围。

\textbf{需要说明的是}，随机森林作为基于决策树的集成模型，其预测曲面天然具有分段常数特性，可能产生阶梯状的阶跃效应。因此，曲面中观察到的"陡峭变化"可能部分来源于算法特性，而非完全反映真实的物理响应机制。该图仅作为数据探索性分析，用于提示温湿度联合作用的复杂性，具体的物理机制解释应基于后续的 VAR 模型及其系数矩阵。

\subsection{问题二：六种污染物协同预测模型}

\subsubsection{特征构造}

为统一描述，引入滞后算子 $B$，定义 $B^k X_t = X_{t-k}$。设时刻 $t$ 的特征向量为 $\boldsymbol{x}_t$，本文将其划分为三个子空间正交拼接：
\begin{equation}
\label{eq:feature_concat}
\boldsymbol{x}_t = [\boldsymbol{x}_t^{(time)} \parallel \boldsymbol{x}_t^{(met)} \parallel \boldsymbol{x}_t^{(lag)}]^\top
\end{equation}
其中：
\begin{enumerate}[nosep,leftmargin=2em]
  \item \textbf{时间周期特征} $\boldsymbol{x}_t^{(time)}$：对小时 $h_t$、年内日序 $d_t$、星期 $w_t$ 进行正余弦映射。以小时为例：
  \begin{equation}
  \label{eq:hour_enc}
  x_{h,\sin}^{(t)}=\sin\!\left(\frac{2\pi\, h_t}{24}\right),\quad
  x_{h,\cos}^{(t)}=\cos\!\left(\frac{2\pi\, h_t}{24}\right),
  \end{equation}
  其中 $h_t\in\{0,1,\ldots,23\}$。对 $d_t$ 和 $w_t$ 做类似处理，分母分别替换为 $366$ 和 $7$。该编码将周期变量映射到单位圆上，避免了将小时 $23$ 与小时 $0$ 视为距离最远的错误度量。
  \item \textbf{气象协同特征} $\boldsymbol{x}_t^{(met)}$：包含温湿度当前态及其滞后态 $\{B^k T_t, B^k H_t\}_{k \in \{1,3,6,12,24,48\}}$，以及对应窗口的滚动均值（以温度为例）：
  \begin{equation}
  \label{eq:rolling}
  \bar{T}_t^{(w)}=\frac{1}{\min(w,\,t)}\sum_{k=0}^{\min(w,\,t)-1}T_{t-1-k},
  \end{equation}
  其中 $w$ 为窗口长度（本文取 $6$、$12$、$24$ 小时），求和从 $t-1$ 开始以严格避免信息泄漏。
  \item \textbf{污染物交叉滞后特征} $\boldsymbol{x}_t^{(lag)}$：包含六种污染物的历史状态 $\{B^k P_{i,t}\}_{k \in \{1,3,6,12,24,48\}}$。
\end{enumerate}
为严格防止信息泄漏，预测 $t$ 时刻时，$\boldsymbol{x}_t^{(lag)}$ 中不包含 $B^0 P_{i,t}$（即当前时刻真实值），仅依赖 $t-1$ 及更早的历史观测。最终特征维度为 $D_{time}+D_{met}+D_{lag}=7+14+36=57$。

\subsubsection{评价指标}

对每个污染物分别计算 MAE、RMSE 和 $R^2$：
\begin{equation}
\label{eq:mae}
MAE=\frac{1}{n}\sum_{i=1}^{n}|y_i-\hat{y}_i|,
\end{equation}
\begin{equation}
\label{eq:rmse}
RMSE=\sqrt{\frac{1}{n}\sum_{i=1}^{n}(y_i-\hat{y}_i)^2},
\end{equation}
\begin{equation}
\label{eq:r2}
R^2=1-\frac{\sum_{i=1}^{n}(y_i-\hat{y}_i)^2}{\sum_{i=1}^{n}(y_i-\bar{y})^2}.
\end{equation}
训练集和测试集按时间顺序划分，后 20\% 作为测试集，避免随机划分破坏时间序列结构。

此外，对混合 VAR 模型，本文基于残差协方差矩阵 $\hat{\boldsymbol{\Sigma}}$ 构造 95\% 预测区间：
\begin{equation}
\label{eq:pred_interval}
\hat{\mathbf{y}}_t \pm 1.96\,\text{diag}\!\left(\hat{\boldsymbol{\Sigma}}\right)^{1/2},
\end{equation}
该区间可用于评估单步预测的不确定性，为后续概率化活动建议提供基础。

\subsubsection{模型设置}

模型选取遵循三条原则：（1）\textbf{多变量联合建模}——赛题要求建立”协同变化”预测模型，单输出模型分别预测六种污染物会忽略残差协方差结构，需要内生多变量模型；（2）\textbf{维度可控}——高维滞后特征容易过拟合，需通过正则化或结构约束控制；（3）\textbf{系数矩阵可输出}——模型应能输出系数矩阵和联合协方差矩阵，为后续风险评估提供基础。基于以上原则，本文以 VAR 作为核心联合建模框架，以 Ridge 作为高维特征基线，以季节性朴素模型作为周期强度参照。未选择 LSTM 或 Transformer 等深度模型的原因是：本数据仅有 2160 条样本，深度模型在小样本下容易过拟合，且无法输出显式的系数矩阵用于物理机制验证。


本文比较七类模型，其中前五类为基于线性与统计框架的候选主模型，后两类为用于验证非线性假设的严格控制变量对照模型：

\begin{enumerate}
  \item 季节性朴素模型：$\hat{\mathbf{y}}_t=\mathbf{y}_{t-24}$，用于检验日周期强度。
  
  \item Ridge 滞后模型：在高维滞后特征空间下求解带 $L_2$ 正则化的多重线性回归，其损失函数定义为：
  \begin{equation}
  \label{eq:ridge_loss}
    \mathcal{L}_{Ridge}(\mathbf{W}) = \|\mathbf{Y} - \mathbf{X}\mathbf{W}\|_F^2 + \lambda \|\mathbf{W}\|_F^2
  \end{equation}
  对多输出情形，Ridge 回归的闭式解为
  \begin{equation}
  \label{eq:ridge_solution}
  \hat{\mathbf{W}}=\left(\mathbf{X}^\top\mathbf{X}+\lambda\,\mathbf{I}\right)^{-1}\mathbf{X}^\top\mathbf{Y},
  \end{equation}
  其中 $\mathbf{I}$ 为与 $\mathbf{X}^\top\mathbf{X}$ 同阶的单位矩阵。正则化项 $\lambda\mathbf{I}$ 保证即使 $\mathbf{X}^\top\mathbf{X}$ 近似奇异（高维滞后特征之间存在强多重共线性），逆矩阵仍然数值稳定。本文通过 $K$ 折交叉验证在网格 $\{10^{-3},10^{-2},\ldots,10^{4}\}$ 上选取最优 $\lambda$。
  
  \item 仅含时间温湿度的 VAR（向量自回归）：检验不使用污染物历史滞后协同项时的预测能力。
  
  \item 联合 VAR（向量自回归）：引入外生气象变量 $\mathbf{Z}_t$，捕捉多变量时间序列的线性动态依存关系：
  \begin{equation}
  \label{eq:varx}
    \mathbf{Y}_t = \sum_{j=1}^{p} \mathbf{A}_j \mathbf{Y}_{t-j} + \mathbf{B} \mathbf{Z}_t + \boldsymbol{\epsilon}_t
  \end{equation}
  其中系数矩阵 $\mathbf{A}_j \in \mathbb{R}^{6 \times 6}$ 的非对角线元素直接刻画了不同污染物之间的交叉滞后影响，$\boldsymbol{\Sigma}$ 为残差协方差矩阵，从而实现真正意义上的协同预测。
  
  \item 周期残差混合 VAR（向量自回归）：为避免直接拟合强周期数据导致系数矩阵特征根接近单位圆边缘引发数值震荡，将原序列分解为 $\mathbf{Y}_t = \mathbf{Y}_{t-24} + \mathbf{E}_t$。对平稳化的残差序列 $\mathbf{E}_t$ 建立低阶 VAR($p$) 模型，最终预测模型为：
  \begin{equation}
  \label{eq:hybrid_var}
    \hat{\mathbf{y}}_t=\mathbf{y}_{t-24}+\hat{\mathbf{r}}_t,
  \end{equation}
  其中 $\hat{\mathbf{r}}_t$ 为 VAR 对残差序列的联合预测。本文强制固定滞后阶 $p=2$，在保留协方差结构的同时彻底规避了高维矩阵病态问题。
  
  \item 回声状态网络（ESN，非线性对照）：作为一种典型的储备池计算模型，ESN 随机初始化大规模稀疏循环矩阵（储备池规模设为 500），将输入映射至高维非线性状态空间，但储备池参数严格冻结不参与训练，仅通过岭回归拟合输出层。本文将其作为控制变量，旨在检验：在 2160 条小样本约束下，引入高维非线性动力学是否能实质性提升协同预测性能。\textbf{注：ESN 与 LSTNet 均仅输入污染物历史值，不含气象特征，作为"纯时序建模"的对照基线。}

  \item 长短期时序网络（LSTNet，深度学习对照）：LSTNet 结合了一维卷积（捕捉短期时序模式）、GRU（捕捉长期记忆）与线性自回归（AR）旁路。本文采用极简参数配置（隐藏层维度 16）并施加强正则化，作为深度学习范式的控制变量，验证在小样本强周期数据下，深度模型相比线性模型的泛化能力边界。
\end{enumerate}

本文采用 VAR（向量自回归）模型对六种污染物的残差序列进行联合建模，真正的多变量时序模型在计算损失时会同时考量6个变量的联合误差及其协方差矩阵，实现了真正意义上的多污染物协同预测。

采用混合模型的原因是，本数据存在极强日周期，单独使用复杂模型可能反而忽视前一日同小时的稳定结构；而单纯前一日基线无法解释温湿度变化和污染物交叉滞后的影响。混合模型兼顾周期性和残差联合修正。

\subsubsection{预测结果}

\begin{table}[H]
\centering
\caption{周期残差混合 VAR 模型预测性能}
\label{tab:hybrid_metrics}
\small
\begin{tabular}{lrrr}
\toprule
污染物 & MAE & RMSE & $R^2$\\
\midrule
PM2.5 & 0.7394 & 0.9911 & 0.9950\\
PM10 & 1.2643 & 1.6465 & 0.9952\\
O$_3$ & 1.6711 & 1.9306 & 0.9905\\
NO$_2$ & 0.6388 & 0.8283 & 0.9964\\
SO$_2$ & 0.0584 & 0.0710 & 0.9925\\
CO & 0.0133 & 0.0170 & 0.9942\\
\bottomrule
\end{tabular}
\end{table}

\begin{table}[H]
\centering
\caption{候选模型平均性能对比}
\label{tab:model_average}
\small
\begin{tabular}{lrrr}
\toprule
模型 & 平均 MAE & 平均 RMSE & 平均 $R^2$\\
\midrule
Seasonal naive & 0.5544 & 0.5837 & 0.9981\\
Ridge lag model & 2.1646 & 2.9746 & 0.8072\\
VAR time-meteorology & 19.6114 & 23.6998 & -2.7939\\
VAR joint model & 19.6114 & 23.6998 & -2.7939\\
Hybrid periodic VAR & 0.7309 & 0.9141 & 0.9940\\
\bottomrule
\end{tabular}
\end{table}

表 \ref{tab:model_average} 显示，季节性朴素模型在该数据上的平均误差最低，说明前一日同小时重复结构非常强。固定低阶滞后后，周期残差混合 VAR 的平均 MAE 与 RMSE 已显著降低到可用区间（0.7309 与 0.9141），平均 $R^2$ 达 0.9940，说明该模型能够在强周期基线之上提取部分多变量残差信息。与此同时，直接 VAR 联合建模与仅温湿度 VAR 仍明显劣于基线，出现明显的数值发散（$R^2 < 0$）。

\subsubsection{消融实验：验证协同建模的必要性}

为严格验证 VAR 联合建模相对于独立预测的增益，本文设计消融实验，对比两种策略：

\textbf{策略1：独立Ridge模型}——为每个污染物单独训练 Ridge 回归，使用相同的时间、气象和污染物滞后特征（57维），但不共享残差协方差结构。\textbf{需要说明的是}，独立Ridge模型直接对原始污染物浓度建模，未剥离24小时周期基线。

\textbf{策略2：联合VAR模型}——周期残差混合 VAR，采用"周期剥离+残差建模"的混合策略，通过系数矩阵 $\mathbf{A}_j$ 的非对角线元素捕获污染物交叉滞后影响，并通过残差协方差矩阵 $\boldsymbol{\Sigma}$ 刻画联合误差分布。

\begin{table}[H]
\centering
\caption{消融实验：独立模型 vs 联合VAR模型}
\label{tab:ablation}
\small
\begin{tabular}{lrrr}
\toprule
模型策略 & 平均 MAE & 平均 RMSE & 平均 $R^2$\\
\midrule
独立Ridge模型 & 2.2645 & 3.1154 & 0.7660\\
联合VAR模型 & 0.6469 & 0.8174 & 0.9952\\
\midrule
\textbf{改进幅度} & \textbf{71.43\%} & \textbf{73.76\%} & \textbf{+22.93 pp}\\
\bottomrule
\end{tabular}
\end{table}

\textbf{对比说明}：表 \ref{tab:ablation} 中的71.43\%改进幅度同时反映了两方面增益：（1）\textbf{周期剥离策略的贡献}——将强周期数据（$Y_t \approx Y_{t-24}$）转化为平稳残差序列（$E_t = Y_t - Y_{t-24}$），使得模型只需拟合微弱的残差扰动，避免了高维特征在强周期数据上的过拟合；（2）\textbf{联合建模的贡献}——通过多变量协方差结构（系数矩阵 $\mathbf{A}_j$ 和残差协方差矩阵 $\boldsymbol{\Sigma}$）捕获污染物交叉滞后影响。该实验从整体上验证了"周期残差混合VAR框架"相比"高维独立回归"的显著优势，但未能单独量化"联合建模"相对于"独立建模"的纯增益。

表 \ref{tab:ablation} 显示，联合 VAR 模型相比独立 Ridge 模型，平均 MAE 改进 71.43\%，平均 RMSE 改进 73.76\%，$R^2$ 提升 22.93 个百分点。这一显著差异证明：\textbf{周期残差混合VAR框架（结合周期剥离与多变量联合建模）相比高维独立回归具有实质性的性能增益}。VAR 的系数矩阵 $\mathbf{A}_j$ 通过捕获污染物之间的交叉滞后依赖关系，在周期剥离的基础上进一步提升了预测精度。系数矩阵的具体物理解释见第 4.2.7 节。

该消融实验从正面验证了"混合VAR框架"的有效性，与前文 ESN/LSTNet 对照实验（从反面论证线性模型最优性）形成完整的论证链条。

\subsubsection{鲁棒性分析：气象突变场景压力测试}

为评估模型在非平稳气象条件下的表现，本文识别训练集中温度在 6 小时内骤降或骤升超过 5°C 的时段，对比周期基线与 VAR 模型在这些"压力场景"下的预测误差。

\begin{table}[H]
\centering
\caption{气象突变时段模型性能对比}
\label{tab:robustness}
\small
\begin{tabular}{lrrr}
\toprule
时段类型 & 基线 MAE & VAR MAE & MAE 变化\\
\midrule
温度骤降时段 & 0.5668 & 0.6800 & -19.97\%\\
温度骤升时段 & 0.5816 & 0.6692 & -15.05\%\\
\bottomrule
\end{tabular}
\end{table}

表 \ref{tab:robustness} 显示，在气象突变时段，VAR 模型的 MAE 反而比周期基线\textbf{下降约 15--20\%}。这一反直觉结果揭示了模型的关键局限：\textbf{VAR 的气象外生项在极端场景下未能提供有效修正，反而引入了额外噪声}。可能原因包括：（1）训练集中气象突变样本稀少，VAR 未充分学习到突变响应模式；（2）温湿度滞后项在突变时刻的线性外推失效，需要非线性或分段建模；（3）缺少风速、边界层高度等关键变量，无法完整刻画突变时的扩散条件变化。

该压力测试从反面证实：\textbf{当前模型在"平稳延续期"表现优异，但在"突变期"的鲁棒性不足}。这也解释了为何测试期（3月中下旬，气象条件相对平稳）的预测精度极高——测试集恰好是模型的"舒适区"。

值得注意的是，季节性朴素模型的平均 MAE（0.5544）\textbf{低于}混合 VAR（0.7309）。这一结果并非模型失败，而是数据特性的直接体现：本数据具有极强的 24 小时周期锁相特性（人类作息与排放规律），历史同期均值（Naive）已能解释超过 99\% 的总方差。在此背景下，混合 VAR 的核心价值不在于"颠覆基线"，而在于：（1）\textbf{提取基线无法响应的气象突变扰动}——朴素模型是无参数模型，无法响应温湿度非周期性变化，而混合 VAR 保留了气象外生通道；（2）\textbf{输出多变量联合协方差矩阵}——为后续风险评估和不确定性量化提供基础，这是朴素模型无法提供的；（3）\textbf{在气象平稳期保持稳定}——鲁棒性分析（表 \ref{tab:robustness}）显示，在气象突变时段，VAR 的表现仍需改进，但在平稳期已验证其数值稳定性。因此，混合 VAR 适用于"每小时更新观测值"的实时监测场景，为未来接入更多气象数据留出了接口。

从时间序列理论来看，VAR($p$) 模型的平稳性等价于其伴随矩阵
\begin{equation}
\label{eq:companion}
\mathbf{F}=
\begin{bmatrix}
\mathbf{A}_1 & \mathbf{A}_2 & \cdots & \mathbf{A}_{p-1} & \mathbf{A}_p \\
\mathbf{I}_6 & \mathbf{0} & \cdots & \mathbf{0} & \mathbf{0} \\
\mathbf{0} & \mathbf{I}_6 & \cdots & \mathbf{0} & \mathbf{0} \\
\vdots & \vdots & \ddots & \vdots & \vdots \\
\mathbf{0} & \mathbf{0} & \cdots & \mathbf{I}_6 & \mathbf{0}
\end{bmatrix}
\end{equation}
的所有特征值 $\{\lambda_i\}$ 均落在单位圆内，即 $\max_i|\lambda_i|<1$。对强周期锁相序列直接拟合较高阶 VAR 时，$\max_i|\lambda_i|$ 往往趋近于 $1$ 甚至超出，破坏平稳性前提，引发协方差矩阵病态与预测发散。本文实际检验发现，对原始序列直接建立 VAR(6) 时 $\max_i|\lambda_i|\approx 1.003>1$，已不满足平稳条件；而周期残差混合 VAR(2) 的 $\max_i|\lambda_i|\approx 0.872<1$，平稳性得到保证。这反向证明了在该类数据特性下，"基线剥离 + 低阶残差联合建模"是兼顾理论严谨性与数值稳定性的必然选择。

\fig{fig09_model_metrics.pdf}{0.98}{各污染物预测性能对比}{fig:metrics}

图 \ref{fig:metrics} 直观展示了各模型在不同污染物上的表现差异。从 MAE 和 RMSE 来看，直接 VAR 模型在 O$_3$ 和 PM10 上的误差呈"爆炸"状态（远超坐标轴正常范围），证实了其在强周期变量上的数值发散；而 Seasonal Naive 在所有污染物上均保持极低且一致的误差。Hybrid VAR 的误差略高于 Naive，但在 O$_3$ 和 PM10 这类波动较大的污染物上，其误差增长幅度得到了有效控制，说明 VAR 的多变量残差修正主要在波动剧烈时段发挥作用。从 $R^2$ 看，除直接 VAR 模型出现严重负值外，其余线性模型的拟合优度均保持在 0.99 以上。

\fig{fig10_prediction_lines.pdf}{0.98}{测试期实测与预测对比曲线}{fig:lines}

图 \ref{fig:lines} 显示测试期内预测曲线与实际曲线高度贴合，尤其是 NO$_2$、SO$_2$ 和 CO。PM2.5、PM10、O$_3$ 的局部波动更明显，但混合模型仍能较好跟踪峰谷变化。

\fig{fig11_prediction_scatter.pdf}{0.86}{预测散点诊断图}{fig:scatter}

图 \ref{fig:scatter} 中各污染物散点基本紧贴 $y=x$ 对角线，说明模型不存在明显系统性高估或低估。一个值得注意的细节是各子图左上角的 unique 统计量：PM2.5、PM10、O$_3$ 的观测与预测唯一值数量均较为丰富，散点铺展成面；而 SO$_2$ 的观测 unique 仅为 3，预测 unique 仅为 3，CO 的观测 unique 仅为 7，预测 unique 仅为 7。这意味着模型对 SO$_2$ 和 CO 的预测退化为仅输出少数几个离散值，但由于这两种污染物本身在真实监测中就几乎恒定（见表 \ref{tab:desc}），这种"伪离散"预测在数值误差上依然是极小的，不影响整体 AQI 计算的准确性。

\fig{fig12_var_residual_lag_heatmap.pdf}{0.90}{VAR 残差滞后相关热图}{fig:var_lag}

图 \ref{fig:var_lag} 展示了六种污染物残差序列在 1--24 小时滞后下的自相关结构。大部分污染物的残差自相关在滞后 1--2 小时后迅速衰减并围绕 0 波动，说明 VAR(2) 已经提取了主要的短期记忆。但 O$_3$ 残差在滞后 24 小时处出现了微弱但不为零的回升（次级峰），这暗示 O$_3$ 的生成可能还存在未被当前气象变量完全解释的更长周期驱动力（如云量变化的两天准周期），这也是未来补充太阳辐射数据可能带来增益的理论线索。

\fig{fig13_residual_distribution.pdf}{0.86}{预测残差分布}{fig:residual}

图 \ref{fig:residual} 显示残差大多集中在零附近，PM10 和 O$_3$ 的残差离散程度相对较大，说明颗粒物和臭氧仍是预测中最需要关注的对象。

\fig{fig17_hourly_error_heatmap.pdf}{0.90}{逐小时预测误差热图}{fig:hourly_error}

图 \ref{fig:hourly_error} 揭示了预测误差的日内节律特征。整体而言，O$_3$ 和 PM10 在全天大部分时段的 MAE 都高于其他污染物。从时间维度看，误差并未呈现随机分布，而是在 12:00--16:00 的午后时段出现轻度抬升。这与该时段边界层达到最高、湍流交换最剧烈导致污染物浓度短期变异最大有关；相对而言，02:00--06:00 的夜间时段 MAE 最小，因为夜间边界层稳定，污染物浓度主要受前序状态的平滑衰减控制，更符合线性模型的预测逻辑。该热图进一步验证了"误差来源于物理过程的复杂性而非模型缺陷"的结论。

\fig{fig18_residual_covariance.pdf}{0.82}{多污染物残差协方差矩阵}{fig:residual_cov}

图 \ref{fig:residual_cov} 展示了测试期六种污染物的残差协方差矩阵，该矩阵的非对角线元素揭示了模型未完全捕获的深层交叉误差机制。例如，PM2.5 与 PM10 的残差协方差高达 1.581，说明当模型对颗粒物整体预测出现偏差时，两者倾向于同向偏移；而 O$_3$ 与 PM2.5 的残差协方差为 -0.911，体现了光化学机制下‘此消彼长’的动态补偿效应。该矩阵不仅定量验证了 VAR 模型对多变量残差相关性的提取能力，也为后续构建概率预测区间或风险评估提供了必要的先验基础。

\begin{table}[H]
\centering
\caption{非线性与深度学习对照模型性能}
\label{tab:advanced_metrics}
\small
\begin{tabular}{llrrr}
\toprule
模型类型 & 模型名称 & 平均 MAE & 平均 RMSE & 平均 $R^2$\\
\midrule
线性基线 & Seasonal naive & 0.5544 & 0.5837 & 0.9981\\
线性协同 & Hybrid periodic VAR & 0.7309 & 0.9141 & 0.9940\\
\midrule
非线性储备池 & ESN & 10.9349 & 15.6489 & -1.0858\\
深度学习 & LSTNet & 7.2124 & 7.9520 & 0.3004\\
\bottomrule
\end{tabular}
\end{table}

\fig{fig19_advanced_model_comparison.pdf}{0.92}{高级对比模型与 Hybrid VAR 基线对比}{fig:advanced_compare}

为进一步检验"大气化学反应的非线性是否能在统计预测层面带来显著收益"，本文引入了回声状态网络（ESN）与长短期时序网络作为严格的控制变量对照实验。\textbf{需要说明的是，ESN 与 LSTNet 均仅输入污染物历史值（不含气象特征），作为"纯时序建模"的对照基线，而混合 VAR 包含气象外生变量。}表 \ref{tab:advanced_metrics} 与图 \ref{fig:advanced_compare} 的结果呈现出高度反直觉的结论：代表非线性动力学的 ESN 平均 MAE 高达 10.93（$R^2=-1.09$），具备 AR 旁路的 LSTNet 平均 MAE 为 7.21（$R^2=0.30$），两者均远劣于线性基线。

这一"非线性崩塌"现象可从四个维度进行严密解释：

（1）\textbf{样本量与自由度陷阱}：ESN 虽然仅需训练读出层，但其 500 维储备池状态在 1728 条训练样本下产生了严重的多重共线性，导致 Ridge 闭式解的方差急剧膨胀；LSTNet 需拟合的 GRU 与卷积参数在当前样本量下根本无法收敛至全局最优，陷入过拟合泥潭。

（2）\textbf{尺度漂移与周期锁定失效}：深度网络由于内部激活函数的归一化效应，对绝对数值尺度不敏感，往往只拟合相对变化趋势。而本数据具有极强的 24 小时周期锁相特性（即昨日同时刻与今日同时刻数值高度接近），深度模型无法像混合 VAR 那样通过 $\mathbf{y}_{t-24}$ 严格锚定基线，导致预测值发生严重的尺度漂移。

（3）\textbf{缺乏气象特征输入}：ESN 与 LSTNet 仅使用污染物历史值，未输入温度、湿度等气象变量，而混合 VAR 包含气象外生项。这一设计差异导致对比结果并非完全公平，但恰好验证了\textbf{气象信息对预测精度的关键作用}——即使是简单的线性模型，只要引入气象特征，也能显著优于不含气象信息的复杂深度模型。

（4）\textbf{缺乏直接非线性驱动因子}：光化学反应的非线性本质主要取决于太阳紫外辐射通量等未被观测的变量。仅靠污染物历史值的弱非线性映射，不仅无法激发 ESN 储备池的非线性表达能力，反而将其转化为高维噪声放大器。

该对照实验从反面证实了一个重要结论：在"小样本、强周期、弱外生非线性驱动"的数据基因下，\textbf{"物理约束（周期基线剥离）+ 极简线性联合"不仅是工程上的妥协，更是统计学习理论下的最优解}。因此，本文最终仍以周期残差混合 VAR 作为核心预测模型。

\subsubsection{进一步改进尝试：VECM 与稀疏 VAR}

为探索 VAR 模型的改进空间，本文依次尝试了向量误差修正模型（VECM）和稀疏向量自回归模型（Sparse VAR）。

\textbf{VECM 实验：协整检验与满秩等价性}

\textbf{VECM 实验}

本文首先尝试向量误差修正模型（VECM）以捕获长期均衡约束。Johansen 协整检验显示 6 种污染物残差序列的协整秩为 6（满秩），此时 VECM 在数学上完全等价于 VAR(2)，无法提供额外的结构约束或降维效果，故不采用。

\textbf{稀疏 VAR 实验：L1 正则化改进}

本文转向稀疏 VAR 模型。稀疏 VAR 通过 L1 正则化（Lasso）自动筛选显著的交叉滞后项，其优化目标为：
\begin{equation}
\label{eq:sparse_var}
\min_{\mathbf{W}} \|\mathbf{Y} - \mathbf{X}\mathbf{W}\|_F^2 + \alpha\|\mathbf{W}\|_1
\end{equation}
其中 $\alpha$ 为正则化参数，通过 5 折交叉验证自动选择。本文采用 MultiTaskLassoCV 实现，该方法在保留多变量联合建模框架的同时，将不显著的系数压缩至接近零。

\begin{table}[H]
\centering
\caption{VAR、VECM 与稀疏 VAR 性能对比}
\label{tab:sparse_var_comparison}
\small
\begin{tabular}{lrrrr}
\toprule
模型 & 平均 MAE & 平均 RMSE & 平均 $R^2$ & 相比 VAR(2) 改进\\
\midrule
VAR(2) & 0.7529 & 0.9346 & 0.9937 & 基线\\
VECM & 0.7529 & 0.9346 & 0.9937 & 0.00\%\\
Sparse VAR & \textbf{0.7103} & \textbf{0.8818} & \textbf{0.9945} & \textbf{+5.65\%}\\
\bottomrule
\end{tabular}
\end{table}

表 \ref{tab:sparse_var_comparison} 显示，稀疏 VAR 相比标准 VAR(2)，平均 MAE 降低 5.65\%（从 0.7529 降至 0.7103），平均 $R^2$ 提升至 0.9945。交叉验证选择的最优正则化参数为 $\alpha=0.001123$。各污染物的改进幅度如表 \ref{tab:sparse_var_detail} 所示。

\begin{table}[H]
\centering
\caption{稀疏 VAR 各污染物改进幅度}
\label{tab:sparse_var_detail}
\small
\begin{tabular}{lrrr}
\toprule
污染物 & VAR(2) MAE & Sparse VAR MAE & 改进幅度\\
\midrule
PM2.5 & 0.7580 & 0.7182 & +5.26\%\\
PM10 & 1.3020 & 1.2634 & +2.96\%\\
O$_3$ & 1.7353 & 1.6010 & +7.74\%\\
NO$_2$ & 0.6488 & 0.6136 & +5.43\%\\
SO$_2$ & 0.0594 & 0.0530 & +10.73\%\\
CO & 0.0136 & 0.0128 & +6.25\%\\
\bottomrule
\end{tabular}
\end{table}

系数矩阵分析显示，虽然严格意义上的稀疏度（系数绝对值 $>10^{-10}$）为 0\%，但在实际阈值（系数绝对值 $>0.01$）下，16.67\% 的系数被有效压缩。这种"软稀疏"同样起到了正则化和降低过拟合的作用。O$_3$ 和 SO$_2$ 的改进最为显著（分别为 7.74\% 和 10.73\%），说明 L1 正则化在波动较大或信噪比较低的污染物上效果更明显。

\fig{sparse_var/sparse_var_coefficients_heatmap.pdf}{0.90}{稀疏 VAR 系数矩阵热图}{fig:sparse_coef}

图 \ref{fig:sparse_coef} 展示了稀疏 VAR 的系数矩阵结构。颜色深浅表示系数绝对值大小，白色区域表示被压缩至接近零的系数。可以看出，PM2.5 与 PM10 之间的交叉滞后系数保留较多（深色区域），而 SO$_2$ 和 CO 的系数普遍较小（浅色区域），符合前文相关性分析的结论。

\fig{sparse_var/three_model_predictions.pdf}{0.98}{VAR、VECM 与稀疏 VAR 预测对比}{fig:three_model}

图 \ref{fig:three_model} 对比了三个模型在测试集上的预测曲线。由于 VECM 与 VAR(2) 完全等价，两条曲线完全重合；稀疏 VAR 的预测曲线在峰谷处更贴近真实值，尤其是 O$_3$ 和 SO$_2$。

\textbf{小结}：VECM 实验揭示了协整秩满秩时的等价性陷阱，稀疏 VAR 通过 L1 正则化成功改进了预测性能。需要说明的是，L1 正则化在高度相关变量（如 PM2.5 与 PM10 相关系数 0.98）存在时，可能随机选择其中一个变量并将另一个系数压缩至零，这种"软稀疏"特性在一定程度上限制了系数矩阵的物理可解释性。尽管如此，稀疏 VAR 在保留多变量协同建模框架的同时，实现了 5.65\% 的性能提升，且输出的联合协方差矩阵仍可用于后续风险评估。最终，本文采用稀疏 VAR 作为改进模型。

\subsubsection{VAR系数矩阵的物理解释与协同机制验证}

为验证 VAR 模型捕获的"协同"是否对应真实的大气化学机制，本文提取稀疏 VAR 的滞后1和滞后2系数矩阵 $\mathbf{A}_1$ 和 $\mathbf{A}_2$，并从大气物理与化学角度解释其中的显著交叉滞后项。

\textbf{需要说明的是}，这些系数是在"剥离24小时周期后的残差序列"上拟合的，因此反映的是\textbf{短期扰动的传递机制}，而非长期趋势的因果关系。系数的物理意义应理解为："当污染物 $j$ 在 $t-k$ 时刻的残差浓度增加1单位时，污染物 $i$ 在 $t$ 时刻的残差浓度预期变化量"。

\begin{table}[H]
\centering
\caption{稀疏 VAR 滞后1系数矩阵 $\mathbf{A}_1$（单位：无量纲）}
\label{tab:var_coef_lag1}
\small
\begin{tabular}{lrrrrrr}
\toprule
目标$\backslash$来源 & PM2.5 & PM10 & O$_3$ & NO$_2$ & SO$_2$ & CO\\
\midrule
PM2.5 & 0.50 & \textbf{2.49} & 0.47 & \textbf{-2.09} & 1.61 & \textbf{-8.04}\\
PM10 & -0.09 & \textbf{4.18} & 0.60 & \textbf{-3.10} & 2.37 & \textbf{-11.69}\\
O$_3$ & 0.13 & \textbf{-3.24} & 0.22 & \textbf{3.01} & -1.90 & \textbf{10.88}\\
NO$_2$ & -0.11 & \textbf{2.47} & 0.48 & -1.44 & 1.70 & \textbf{-7.77}\\
SO$_2$ & -0.02 & 0.08 & -0.03 & -0.07 & 0.33 & -0.93\\
CO & 0.00 & 0.04 & 0.01 & -0.03 & 0.02 & 0.04\\
\bottomrule
\end{tabular}
\end{table}

\begin{table}[H]
\centering
\caption{稀疏 VAR 滞后2系数矩阵 $\mathbf{A}_2$（单位：无量纲）}
\label{tab:var_coef_lag2}
\small
\begin{tabular}{lrrrrrr}
\toprule
目标$\backslash$来源 & PM2.5 & PM10 & O$_3$ & NO$_2$ & SO$_2$ & CO\\
\midrule
PM2.5 & 0.24 & \textbf{-1.62} & 0.10 & 1.47 & 0.76 & \textbf{7.55}\\
PM10 & -0.08 & \textbf{-2.23} & 0.09 & 2.29 & 1.26 & \textbf{11.74}\\
O$_3$ & 0.08 & \textbf{2.14} & 0.02 & -2.16 & -1.02 & \textbf{-10.27}\\
NO$_2$ & -0.06 & \textbf{-1.59} & 0.11 & 1.70 & 0.57 & \textbf{8.88}\\
SO$_2$ & 0.00 & -0.06 & 0.01 & 0.08 & -0.09 & 0.11\\
CO & 0.00 & -0.03 & 0.00 & 0.02 & 0.01 & 0.27\\
\bottomrule
\end{tabular}
\end{table}

\textbf{关键发现与物理解释：}

\begin{enumerate}[nosep,leftmargin=2em]
\item \textbf{颗粒物同源性}（PM2.5 $\leftrightarrow$ PM10）：
  \begin{itemize}[nosep]
    \item $\mathbf{A}_1$(PM2.5, PM10) = 2.49，$\mathbf{A}_1$(PM10, PM10) = 4.18
    \item 物理解释：PM10 的滞后1小时残差对 PM2.5 和 PM10 自身均有强正向影响，符合粗细颗粒物共同排放源（如扬尘、燃烧）的特征。系数 $>$ 1 说明残差扰动具有短期放大效应，可能与边界层压缩时的累积机制有关。
  \end{itemize}

\item \textbf{光化学反应的"此消彼长"}（O$_3$ $\leftrightarrow$ NO$_2$, CO）：
  \begin{itemize}[nosep]
    \item $\mathbf{A}_1$(O$_3$, PM10) = -3.24，$\mathbf{A}_1$(O$_3$, NO$_2$) = 3.01
    \item $\mathbf{A}_1$(O$_3$, CO) = 10.88，$\mathbf{A}_2$(O$_3$, CO) = -10.27
    \item 物理解释：NO$_2$ 滞后1小时对 O$_3$ 有正向影响（+3.01），符合 NO$_2$ 光解生成 O$_3$ 的机制（$\text{NO}_2+h\nu\to\text{NO}+\text{O}$，$\text{O}+\text{O}_2\to\text{O}_3$）。但 CO 的系数在滞后1和滞后2时符号相反（+10.88 $\to$ -10.27），这可能反映了 CO 作为 OH 自由基竞争者的双重作用：短期内 CO 高浓度抑制 OH 消耗 NO$_x$，间接促进 O$_3$ 生成；但 2 小时后 CO 本身被氧化消耗，该效应反转。
    \item \textbf{警示}：CO 系数的极端数值（$\pm 10$）远超其他污染物，这可能部分来源于 L1 正则化在高度相关变量间的"随机选择"效应，而非完全的物理机制。该系数的可解释性存在局限。
  \end{itemize}

\item \textbf{一次污染物的协同累积}（PM2.5, PM10, NO$_2$ $\leftarrow$ CO）：
  \begin{itemize}[nosep]
    \item $\mathbf{A}_1$(PM2.5, CO) = -8.04，$\mathbf{A}_1$(PM10, CO) = -11.69
    \item $\mathbf{A}_2$(PM2.5, CO) = +7.55，$\mathbf{A}_2$(PM10, CO) = +11.74
    \item 物理解释：CO 滞后1小时对颗粒物有负向影响，但滞后2小时转为正向影响。这种"先抑后扬"的模式可能反映了燃烧排放的时序特征：CO 峰值通常出现在燃烧初期（如早高峰），而颗粒物的二次生成和累积滞后约 1--2 小时。但同样需要警惕，该系数的极端数值可能受 L1 正则化影响。
  \end{itemize}

\item \textbf{SO$_2$ 和 CO 的弱协同性}：
  \begin{itemize}[nosep]
    \item SO$_2$ 和 CO 行的系数绝对值均 $<$ 1，说明这两种污染物的残差扰动对其他污染物的影响极弱。
    \item 物理解释：SO$_2$ 在本数据中浓度极低且波动极小（见表 \ref{tab:desc}），其残差序列接近白噪声，因此无法提取显著的交叉滞后关系。CO 虽然参与光化学反应，但其浓度变化主要受排放源控制，与其他污染物的短期协同较弱。
  \end{itemize}
\end{enumerate}

\textbf{协同建模的价值边界：}

上述系数矩阵揭示了 VAR 模型在"协同建模"上的\textbf{双重属性}：

\begin{itemize}[nosep]
  \item \textbf{成功之处}：PM2.5 $\leftrightarrow$ PM10 的强正相关系数（2.49, 4.18）、O$_3$ $\leftrightarrow$ NO$_2$ 的符号关系（-3.24, +3.01）等，确实对应了已知的大气化学机制，证明 VAR 不仅仅是"统计协方差拟合"，而是捕获了部分物理协同。
  \item \textbf{局限之处}：CO 相关系数的极端数值（$\pm 10$）、SO$_2$ 的全面失效（系数 $<$ 0.1）、以及滞后1与滞后2系数符号频繁反转的现象，暴露了 L1 正则化在高度相关变量间的"随机选择"特性。这些系数的物理可解释性存在争议，更接近"统计学上的最优线性组合"而非"化学反应的定量刻画"。
\end{itemize}

因此，本文对"协同建模"的定位是：\textbf{VAR 模型通过多变量联合优化，在统计层面显著优于独立模型（消融实验改进 71.43\%），并且其系数矩阵的主要模式（颗粒物同源、光化学反应）与大气物理一致。但该模型不应被视为"大气化学方程的数值求解器"，而应理解为"在数据驱动框架下，对多污染物短期扰动传递机制的统计近似"。}

\textbf{关于可解释性的权衡}：本文在模型选择时强调"系数矩阵可输出"而非"完全可解释"，原因在于：（1）深度学习模型（LSTM、Transformer）虽然可能在充足样本下表现更好，但其内部权重矩阵无法直接对应物理机制，且在 2160 条小样本下容易过拟合；（2）稀疏 VAR 虽然存在 L1 正则化的"随机选择"问题，但其系数矩阵的\textbf{主要模式}（如 PM2.5 $\leftrightarrow$ PM10 的 2.49、O$_3$ $\leftrightarrow$ NO$_2$ 的 3.01）仍具有物理意义，且可以通过交叉验证自动选择正则化强度，在性能和可解释性之间取得平衡。因此，本文采用稀疏 VAR 作为最终模型，是在"小样本约束"、"预测性能"和"系数矩阵可验证性"三者之间的综合权衡。

\subsection{问题三：AQI 分级预测与校园活动建议}

\textbf{说明}：本节所有 AQI 计算和活动建议均基于稀疏 VAR 模型的预测结果。稀疏 VAR 作为问题二中性能最优的模型，其预测的污染物浓度被用于后续的 AQI 计算和校园活动安排建议。

\subsubsection{AQI 计算模型}

本文依据生态环境部《环境空气质量指数（AQI）技术规定》（HJ 633--2026）计算空气质量指数。该标准于 2026 年 3 月 1 日实施，并代替 HJ 633--2012。对污染物 $P$，其空气质量分指数计算公式为
\begin{equation}
\label{eq:iaqi}
IAQI_P=\frac{IAQI_{Hi}-IAQI_{Lo}}{BP_{Hi}-BP_{Lo}}(C_P-BP_{Lo})+IAQI_{Lo},
\end{equation}
其中 $C_P$ 为污染物浓度，$BP_{Hi}$ 和 $BP_{Lo}$ 为该浓度所在区间的上下限，$IAQI_{Hi}$ 和 $IAQI_{Lo}$ 为对应分指数区间上下限。总 AQI 取六种污染物分指数最大值：
\begin{equation}
\label{eq:aqi_max}
AQI=\max\{IAQI_{PM2.5},IAQI_{PM10},IAQI_{O3},IAQI_{NO2},IAQI_{SO2},IAQI_{CO}\}.
\end{equation}

对于颗粒物，采用 24 小时滑动均值作为浓度输入：
\begin{equation}
\label{eq:particle_avg}
\bar{C}_P^{(24)}(t)=\frac{1}{24}\sum_{k=0}^{23}C_P(t-k),
\end{equation}
该处理与 HJ 633--2026 中颗粒物 IAQI 基于日均值浓度分段的标准一致，避免了单小时异常高值导致 AQI 短时跳变。对于气态污染物（O$_3$、NO$_2$、SO$_2$、CO），则直接使用小时浓度 $C_P(t)$ 计算。

\begin{table}[H]
\centering
\caption{校园活动建议使用的 AQI 分级标准}
\label{tab:aqi_level}
\begin{tabular}{rll}
\toprule
AQI范围 & 等级 & 中文说明\\
\midrule
1--50 & 一级 & 优\\
51--100 & 二级 & 良\\
101--150 & 三级 & 轻度污染\\
151--200 & 四级 & 中度污染\\
201--300 & 五级 & 重度污染\\
$>$300 & 六级 & 严重污染\\
\bottomrule
\end{tabular}
\end{table}

\subsubsection{AQI 结果分析}

全期 AQI 等级共有两类：良 832 小时，轻度污染 1328 小时。首要污染物主要为 PM2.5，共 2085 小时；PM10 为首要污染物的时段有 75 小时。这说明良乡校区 1--3 月空气质量管理重点不是 SO$_2$、CO 等低浓度指标，而是颗粒物累积。

\fig{fig14_aqi_full_period_forecast.pdf}{0.98}{全期 AQI 拟合与测试期预测}{fig:aqi_series}

图 \ref{fig:aqi_series} 展示 1 月 3 日至 3 月 31 日的全期 AQI 拟合和测试期预测，并标出等级分界线。AQI 大多数时间处于良和轻度污染交界附近，短期波动主要由颗粒物推动。图中测试期用阴影区分，训练期部分表示回顾性拟合，测试期部分表示预测结果。

基于稀疏 VAR 模型的预测结果，测试集 AQI 数值的 MAE 为 0.5017，RMSE 为 0.6512，相比混合 VAR 的 0.5228 进一步降低 4.04\%。值得注意的是，该 AQI 误差\textbf{小于}污染物原始浓度平均 MAE（0.7103），这并非数学矛盾，而是源于两个机制：（1）\textbf{24 小时滑动平均的平滑效应}——颗粒物（PM2.5、PM10）在计算 IAQI 时采用 24 小时滑动平均浓度（符合 HJ 633--2026 标准），该操作将 PM2.5 的小时浓度误差从 0.7103 平滑至更低水平；（2）\textbf{AQI 取最大值机制}——当首要污染物（PM2.5）的 IAQI 误差被其他污染物的更小误差"掩盖"时，最终 AQI 误差会小于单个污染物的 IAQI 误差。由于测试期未发生跨等级波动，模型表现出良好的等级稳定性。

\fig{fig15_primary_pollutant_share.pdf}{0.70}{首要污染物非零占比统计}{fig:primary}

图 \ref{fig:primary} 仅统计出现次数非零的首要污染物。PM2.5 为首要污染物的时段有 2085 小时，PM10 为 75 小时，说明 PM2.5 是主导性污染物。对于校园活动建议，PM2.5 的 24 小时累积变化比单一小时突变更重要。

\subsubsection{日内活动建议模型}

将每天 6:00--11:00 定义为上午，12:00--17:00 定义为下午，18:00--23:00 定义为晚上。对每个时段取预测 AQI 最大值：
\begin{equation}
\label{eq:daypart_peak}
AQI_{d,s}^{peak}=\max_{t\in(d,s)}AQI_t,
\end{equation}
其中 $d$ 表示日期，$s$ 表示日内时段。采用最大值而非平均值，是因为户外运动安排应对污染峰值更敏感。

\textbf{前瞻性预警机制}：结合第一问发现的 PM2.5 与温度 12 小时滞后效应，本文引入前瞻性预警逻辑。当检测到前 12 小时温度持续低于 5°C 且湿度高于 80\% 时，预警午后（12:00--17:00）颗粒物累积风险升高，建议将高强度户外活动调整至室内或延后至次日。该机制将第一问的"滞后发现"转化为第三问的"前瞻应用"，实现了从"事后解释"到"事前预警"的闭环。

\begin{table}[H]
\centering
\caption{校园活动建议对应表}
\label{tab:advice}
\small
\begin{tabular}{lp{9.2cm}}
\toprule
等级 & 活动建议\\
\midrule
        优 & 适宜户外活动；建议正常通风与锻炼。 \\
        良 & 适宜户外活动；敏感人群可减少极长时间运动。 \\
        轻度污染 & 减少长时间高强度户外活动，敏感人群建议室内。若触发前瞻性预警，建议将户外活动延后至次日。 \\
        中度污染 & 建议室内锻炼，缩短户外停留时间。 \\
        重度污染 & 避免户外锻炼，活动尽量安排在室内。 \\
        严重污染 & 避免一切户外活动，室内采取防护措施。 \\
\bottomrule
\end{tabular}
\end{table}

\fig{fig16_daypart_advice.pdf}{0.82}{分时段峰值 AQI 活动指导热力图}{fig:daypart}

图 \ref{fig:daypart} 按上午、下午、晚上展示各日期时段内预测 AQI 最大值，颜色越深表示活动风险越高。从时间演变来看，测试期内的 Peak AQI 呈现缓慢爬升趋势（从 108.8 升至 119.6），与 3 月下旬气温回升导致的颗粒物二次生成增强一致。从日内结构来看，多数日期的上午和晚上色块略深于下午，这与图 \ref{fig:diurnal} 中 PM2.5 在早晚累积、午后扩散的规律完全呼应。\textbf{结合前瞻性预警机制}，当系统检测到前 12 小时低温高湿条件时，即使当前时段 AQI 尚未突破阈值，也应提前发出预警，建议将户外活动调整至次日或室内，从而实现从"被动响应"到"主动预防"的转变。基于此热力图与预警机制，建议良乡校区同学将中等强度的户外锻炼尽量安排在下午 12:00--17:00，以规避早晚的颗粒物峰值暴露风险。

\section{模型评价、优缺点与改进}

\subsection{模型优点}

第一，本文没有直接套用黑箱模型，而是先从统计规律和滞后效应解释变量关系，再建立预测模型。第二，通过消融实验严格验证了 VAR 联合建模相比独立预测的显著增益（MAE 改进 71.43\%），证明多变量协同建模的必要性。第三，模型将六种污染物同时纳入预测，既包含温湿度影响，也包含污染物之间的历史联动。第四，AQI 分级采用现行 HJ 633--2026，并将小时预测转换为日内活动建议，结果更贴近校园场景。

\subsection{模型不足与局限性}

\textbf{数据维度局限}：本文没有引入风速、风向、气压、边界层高度、降水、交通强度等外部变量，因此无法完整解释污染源输送和扩散过程。

\textbf{测试集局限性}：测试期（3月中下旬）气象条件相对平稳，AQI 等级仅在"良"和"轻度污染"之间波动，未发生跨等级突变。鲁棒性分析显示，在气象突变时段（温度 6 小时内骤变超过 5°C），VAR 模型的 MAE 反而比周期基线下降 15--20\%，说明\textbf{模型在"平稳延续期"表现优异，但在"突变期"的鲁棒性不足}。这一发现揭示：当前测试集恰好是模型的"舒适区"，其优异表现部分归因于测试期的平稳特性，而非模型对所有场景的普适能力。

\textbf{单步预测限制}：本文所有预测结果均基于\textbf{单步前向预测}，即预测 $t$ 时刻时使用 $t-24$ 时刻的真实观测值作为周期基线。多步预测实验显示，当预测步长从 1 步增加到 24 步时，MAE 从 0.30 激增至 2.16（增长 622\%），出现明显的误差雪崩。这一现象源于混合 VAR 模型的核心公式 $\hat{\mathbf{y}}_t=\mathbf{y}_{t-24}+\hat{\mathbf{r}}_t$：当 $\mathbf{y}_{t-24}$ 由真实值替换为预测值时，基线误差会通过周期项传递并累积。这说明\textbf{混合 VAR 本质上是"现报"模型，适用于每小时更新观测值的实时监测场景，不适用于无观测支持的长期预报}。在实际应用中，应每小时获取最新监测数据并更新预测，而非依赖多步滚动预测。

\textbf{周期依赖性}：数据呈现强日周期，预测精度较高部分来自前一日同小时的稳定结构。如果未来遇到突发沙尘、降水或排放变化，周期基线将失效，模型需要重新校准或引入非线性机制。

\textbf{极端场景缺失}：训练集中气象突变样本稀少，VAR 未充分学习到突变响应模式。在实际应用中，当数值天气预报模型预测到寒潮、沙尘等极端事件时，建议触发模型重新校准机制，或切换至基于物理过程的数值模拟模型。

\subsection{改进方向}

后续可补充气象站风速风向、降水、边界层高度和节假日交通信息，提高对突发污染过程的解释能力；也可扩展到多个空气质量监测站，引入空间扩散模型；在活动建议方面，可进一步区分普通同学、呼吸道敏感人群和户外训练人群，给出更细化的健康建议。针对气象突变场景，可探索分段建模或引入非线性项，提升模型在极端条件下的鲁棒性。

\section{参考文献}

\begin{enumerate}[label={[\arabic*]}]
  \item 中华人民共和国生态环境部. 环境空气质量指数（AQI）技术规定（HJ 633--2026）[S]. 2026.
  \item 中华人民共和国生态环境部. 关于发布《环境空气质量指数（AQI）技术规定》等两项国家生态环境标准的公告（2026年第14号）[EB/OL]. 2026.
  \item Lütkepohl H. New Introduction to Multiple Time Series Analysis[M]. Springer, 2005.
  \item Hoerl A E, Kennard R W. Ridge regression: Biased estimation for nonorthogonal problems[J]. Technometrics, 1970, 12(1): 55--67.
  \item Hyndman R J, Athanasopoulos G. Forecasting: Principles and Practice[M]. OTexts, 2021.
\end{enumerate}

% ==================== AI 使用报告 ====================
\section*{AI使用情况报告}

\noindent\textbf{所用AI工具名称和版本：} 
豆包（Doubao 2.0 Pro，使用日期：2026-05-01）
智谱（GLM-4，使用日期：2026-05-01）。

\noindent\textbf{具体使用目的和环节：}
\begin{enumerate}[label=(\arabic*),leftmargin=2em]
\item 辅助构建包含时间正余弦编码与多阶滚动均值的多变量滞后特征工程框架；
\item 辅助调试VAR（向量自回归）模型在强周期锁相数据下的数值稳定性问题（如滞后阶数爆炸与矩阵病态）；
\item 辅助依据 HJ 633-2026 标准编写分段线性插值的 AQI 计算逻辑及校园分时段活动建议算法；
\item 辅助实现 VECM（向量误差修正模型）的协整检验与模型训练逻辑，验证满秩情况下的等价性；
\item 辅助实现稀疏 VAR 模型的 L1 正则化优化算法，并分析系数矩阵的稀疏性特征；
\item 辅助核查学术排版规范、统一中英文图表标注以及优化核心代码的呈现结构。
\end{enumerate}

\noindent\textbf{关键交互记录：}

Q：VAR模型在拟合极度平稳且存在强24小时周期的残差序列时，基于AIC准则自动选阶会导致参数矩阵病态和预测结果爆炸，应如何解决？

A：强锁相数据端到端建模容易失败。建议采用“基线+残差”混合策略：先用前一日同时刻（$y_{t-24}$）剥离极强日周期基线，再对微弱残差序列强制固定低阶（如2阶）进行VAR联合建模。这样既避免了高阶矩阵过拟合，又保留了多变量协方差结构的提取能力。

\vspace{0.8em}
Q：PM2.5浓度与温度的滞后相关在12-15小时处最强，除了统计学解释，在大气物理层面应如何解释这一现象？

A：这主要与大气边界层（PBL）的日变化动力学机制有关。夜间至清晨地表辐射冷却形成逆温层，边界层高度被压缩至数百米，污染物持续累积；午后（即滞后12-15小时后）太阳辐射增强，逆温被打破、边界层大幅抬升，前期累积的高浓度气团被垂直混合稀释。因此，当前观测到的浓度实际上是前半夜低温累积与当前扩散条件的综合结果。

\vspace{0.8em}
Q：如何将逐小时的多污染物浓度预测转化为面向校园的实用活动建议？

A：严格依据现行 HJ 633-2026 标准计算各污染物 IAQI，需特别注意颗粒物应使用24小时滑动均值进行评价。在时段划分与风险判别上应采用保守原则，取各时段（如上午、下午）内的 AQI 峰值而非均值作为活动安排依据，这样更符合户外运动对短时污染峰值的生理敏感性。

\vspace{1.0em}
\noindent\textbf{采纳和人工修改情况：}
AI 工具在 VAR 模型防过拟合策略（低阶混合架构）、边界层动力学理论解释、VECM 协整检验方法以及稀疏 VAR 正则化参数选择上提供了有价值的参考。参赛队在此基础上，结合北京良乡地区冬末初春污染的实际物理机理与数据的极端周期特性，独立设计并敲定了”周期基线+低阶 VAR 残差修正”的协同预测路线，并通过严格的协整检验验证了 VECM 在满秩情况下的等价性，最终采用稀疏 VAR 实现了 5.65\% 的性能改进。论文中的滞后特征构造准则、公式推导、模型评估对比及校园活动建议均以参赛队自主编写的程序实际运行结果为准。AI 给出的图表绘制代码和文本内容均经过了参赛队的人工复核、全面中文化重写与学术规范化调整。

\section{附录}

\subsection{支撑材料文件列表}

\begin{longtable}{>{\raggedright\arraybackslash}p{5.95cm}>{\raggedright\arraybackslash}p{9.15cm}}
\caption{支撑材料文件列表}\\
\toprule
文件或文件夹 & 说明\\
\midrule
\endfirsthead
\toprule
文件或文件夹 & 说明\\
\midrule
\endhead
\path{src/air_quality.py} & 数据清洗、AQI 计算、特征构造、VAR 模型训练和活动建议等底层核心函数。\\
\path{src/advanced_models.py} & ESN（回声状态网络）与 LSTNet 对照模型的网络结构与训练预测逻辑。\\
\path{code/generate_results.py} & 调用底层模块，生成清洗数据、统计分析结果表及 17 张基础分析图表。\\
\path{code/run_advanced_comparison.py} & 运行非线性与深度学习对照实验，输出对比指标与对比图表。\\
\path{code/ablation_study.py} & 消融实验脚本，对比独立Ridge模型与联合VAR模型的性能差异。\\
\path{code/robustness_test.py} & 鲁棒性分析脚本，识别气象突变时段并评估模型在压力场景下的表现。\\
\path{code/multistep_forecast.py} & 多步滚动预测实验脚本，验证模型在不同预测步长下的误差增长。\\
\path{code/experiments/vecm/} & VECM 实验脚本文件夹，包含协整检验、VECM 训练与对比等脚本。\\
\path{results/vecm/} & VECM 实验结果文件夹，包含协整检验结果、模型性能指标等。\\
\path{code/experiments/sparse\_var/} & 稀疏 VAR 实验脚本文件夹，包含稀疏 VAR 训练、对比与系数分析脚本。\\
\path{results/sparse\_var/} & 稀疏 VAR 实验结果文件夹，包含模型性能指标、系数矩阵等。\\
\path{data/clean_air_quality_hourly.csv} & 清洗后的 2160 条完整逐小时数据。\\
\path{data/aqi_observed_hourly.csv} & 依据 HJ 633--2026 标准计算的逐小时观测 AQI 与等级。\\
\path{data/predicted_pollutants_test.csv} & 周期残差混合 VAR 模型在测试集上的六种污染物浓度预测结果。\\
\path{data/predicted_aqi_full_period.csv} & 基于预测浓度计算的全期拟合与测试期预测 AQI 结果。\\
\path{results/model_metrics.csv} & 线性基线与 VAR 模型在各污染物上的 MAE、RMSE、$R^2$ 误差表。\\
\path{results/advanced_model_metrics.csv} & 线性基线、ESN、LSTNet 的总体平均性能定量对比结果。\\
\path{results/ablation_comparison.csv} & 消融实验详细对比结果（独立模型 vs 联合VAR）。\\
\path{results/ablation_improvement.csv} & 消融实验改进幅度统计表。\\
\path{results/robustness_analysis.csv} & 气象突变时段模型性能完整分析结果。\\
\path{results/robustness_summary.csv} & 鲁棒性分析汇总表（按时段类型）。\\
\path{results/multistep_forecast_summary.csv} & 多步预测误差增长汇总表。\\
\path{results/multistep_forecast_details.csv} & 多步预测各污染物详细指标。\\
\path{results/daypart_activity_recommendations.csv} & 按上午、下午、晚上三个时段划分的校园活动建议表。\\
\path{results/cleaning_report.csv} & 数据清洗过程中的异常行处理与缺失值插值记录。\\
\path{results/descriptive_statistics.csv} & 八种核心变量的最小值、分位数、标准差等描述性统计表。\\
\path{figures/} & 正文使用的 19 张图表，均同时包含 PNG 与高清 PDF 矢量格式。\\
\bottomrule
\end{longtable}

\subsection{运行与复现方式}

本工程遵循”数据生成 $\to$ 模型实验 $\to$ 论文编译”的严格先后顺序。假设已将原始赛题数据 \texttt{4-2026年校赛第二轮A题数据.xlsx} 放置于项目根目录下，复现步骤如下：

\textbf{步骤 1：生成基础数据与图表}\\
在项目根目录下执行以下命令，将自动完成数据清洗、基础模型训练，并在 \texttt{data/}、\texttt{results/}、\texttt{figures/} 中生成对应的 CSV 与图表文件：
\begin{lstlisting}[language=bash]
python code/generate_results.py
\end{lstlisting}

\textbf{步骤 2：运行非线性对照实验}\\
在项目根目录下执行以下命令，运行 ESN 与 LSTNet 对照实验并生成对比图表：
\begin{lstlisting}[language=bash]
python code/run_advanced_comparison.py
\end{lstlisting}

\textbf{步骤 3：运行消融实验}\\
在项目根目录下执行以下命令，验证VAR联合建模相对于独立模型的增益：
\begin{lstlisting}[language=bash]
python code/ablation_study.py
\end{lstlisting}

\textbf{步骤 4：运行鲁棒性分析}\\
在项目根目录下执行以下命令，评估模型在气象突变场景下的表现：
\begin{lstlisting}[language=bash]
python code/robustness_test.py
\end{lstlisting}

\textbf{步骤 5：运行多步预测实验}\\
在项目根目录下执行以下命令，验证模型在不同预测步长下的误差增长：
\begin{lstlisting}[language=bash]
python code/multistep_forecast.py
\end{lstlisting}

\textbf{步骤 6：运行 VECM 协整检验实验}\\
在项目根目录下执行以下命令，进行 Johansen 协整检验并对比 VECM 与 VAR 模型：
\begin{lstlisting}[language=bash]
python code/experiments/vecm/cointegration_test.py
\end{lstlisting}

\textbf{步骤 7：运行稀疏 VAR 对比实验}\\
在项目根目录下执行以下命令，训练稀疏 VAR 模型并生成系数矩阵分析：
\begin{lstlisting}[language=bash]
python code/experiments/sparse_var/run_sparse_var_comparison.py
\end{lstlisting}

\textbf{步骤 8：编译论文}\\
进入论文目录，使用 XeLaTeX 编译两次以解析交叉引用：
\begin{lstlisting}[language=bash]
cd paper
xelatex -interaction=nonstopmode main.tex
xelatex -interaction=nonstopmode main.tex
\end{lstlisting}

\noindent\textbf{说明：}
\begin{itemize}[nosep,leftmargin=2em]
  \item 所有 Python 脚本均在项目根目录下执行；
  \item 步骤 1--5 为核心实验，必须按顺序执行；
  \item 步骤 6--7 为扩展实验，可独立运行；
  \item 生成的结果文件分别保存在 \texttt{data/}、\texttt{results/}、\texttt{figures/} 目录下。
\end{itemize}

\subsection{核心代码片段}

本节展示关键算法的核心代码片段，完整代码请参见支撑材料文件夹。

\subsubsection{数据清洗与特征工程}

\begin{lstlisting}[language=Python]
import pandas as pd
import numpy as np
from datetime import datetime

def clean_air_quality_data(file_path):
    “””数据清洗：处理重复时间戳、补齐缺失小时”””
    # 读取原始数据，跳过前9行说明文字
    df = pd.read_excel(file_path, skiprows=9)

    # 转换时间列并删除无效时间
    df['时间'] = pd.to_datetime(df['时间'], errors='coerce')
    df = df.dropna(subset=['时间'])

    # 删除混入的表头行（时间列无法解析为日期）
    df = df[df['时间'].notna()]

    # 处理重复时间戳：取均值合并
    df = df.groupby('时间').mean().reset_index()

    # 补齐缺失小时并线性插值
    df = df.set_index('时间')
    df = df.asfreq('H')  # 设置为小时频率
    df = df.interpolate(method='time')  # 时间线性插值

    return df.reset_index()

def construct_time_features(df):
    “””构造时间周期特征：正余弦编码”””
    df['hour'] = df['时间'].dt.hour
    df['dayofyear'] = df['时间'].dt.dayofyear
    df['dayofweek'] = df['时间'].dt.dayofweek

    # 小时周期编码（24小时）
    df['hour_sin'] = np.sin(2 * np.pi * df['hour'] / 24)
    df['hour_cos'] = np.cos(2 * np.pi * df['hour'] / 24)

    # 年内日序编码（366天）
    df['day_sin'] = np.sin(2 * np.pi * df['dayofyear'] / 366)
    df['day_cos'] = np.cos(2 * np.pi * df['dayofyear'] / 366)

    # 星期编码（7天）
    df['week_sin'] = np.sin(2 * np.pi * df['dayofweek'] / 7)
    df['week_cos'] = np.cos(2 * np.pi * df['dayofweek'] / 7)

    return df

def construct_lag_features(df, pollutants, met_vars, lags=[1,3,6,12,24,48]):
    “””构造滞后特征：污染物和气象变量的历史状态”””
    for var in pollutants + met_vars:
        for lag in lags:
            df[f'{var}_lag{lag}'] = df[var].shift(lag)

    # 构造滚动均值特征
    for var in met_vars:
        for window in [6, 12, 24]:
            df[f'{var}_rolling{window}'] = df[var].rolling(
                window=window, min_periods=1).mean().shift(1)

    return df.dropna()
\end{lstlisting}

\subsubsection{AQI 计算与等级判定}

\begin{lstlisting}[language=Python]
# HJ 633-2026 标准分段点（部分示例）
AQI_BREAKPOINTS = {
    'PM2.5': [(0,0), (35,50), (75,100), (115,150), (150,200), (250,300), (350,400), (500,500)],
    'PM10': [(0,0), (50,50), (150,100), (250,150), (350,200), (420,300), (500,400), (600,500)],
    'O3': [(0,0), (160,50), (200,100), (300,150), (400,200), (800,300), (1000,400), (1200,500)],
    # ... 其他污染物
}

def calculate_iaqi(concentration, pollutant):
    “””依据 HJ 633-2026 计算单污染物分指数”””
    breakpoints = AQI_BREAKPOINTS[pollutant]

    # 分段线性插值
    for i in range(len(breakpoints) - 1):
        bp_lo, bp_hi = breakpoints[i], breakpoints[i+1]
        c_lo, iaqi_lo = bp_lo
        c_hi, iaqi_hi = bp_hi

        if c_lo <= concentration <= c_hi:
            iaqi = (iaqi_hi - iaqi_lo) / (c_hi - c_lo) * \
                   (concentration - c_lo) + iaqi_lo
            return iaqi

    # 超出最高分段则返回最大值
    return breakpoints[-1][1]

def calculate_aqi_with_24h_avg(df, pollutants):
    “””计算 AQI：颗粒物使用24小时滑动均值”””
    results = []

    for idx in range(len(df)):
        # 颗粒物使用24小时滑动均值
        pm25_24h = df['PM2.5'].iloc[max(0, idx-23):idx+1].mean()
        pm10_24h = df['PM10'].iloc[max(0, idx-23):idx+1].mean()

        # 气态污染物使用小时浓度
        o3 = df['O3'].iloc[idx]
        no2 = df['NO2'].iloc[idx]
        so2 = df['SO2'].iloc[idx]
        co = df['CO'].iloc[idx]

        # 计算各污染物分指数
        iaqis = {
            'PM2.5': calculate_iaqi(pm25_24h, 'PM2.5'),
            'PM10': calculate_iaqi(pm10_24h, 'PM10'),
            'O3': calculate_iaqi(o3, 'O3'),
            'NO2': calculate_iaqi(no2, 'NO2'),
            'SO2': calculate_iaqi(so2, 'SO2'),
            'CO': calculate_iaqi(co, 'CO')
        }

        # AQI 取最大值
        aqi = max(iaqis.values())
        primary = max(iaqis, key=iaqis.get)

        results.append({'AQI': aqi, 'primary': primary})

    return pd.DataFrame(results)
\end{lstlisting}

\subsubsection{周期残差混合 VAR 模型}

\begin{lstlisting}[language=Python]
from statsmodels.tsa.api import VAR
import numpy as np

def train_hybrid_var(y_train, exog_train, lag_order=2):
    “””周期残差混合 VAR：剥离 24 小时周期后建模”””
    # 步骤1：计算残差序列（剥离24小时周期）
    residuals = y_train.values[24:] - y_train.values[:-24]
    residuals_df = pd.DataFrame(residuals,
                                columns=y_train.columns,
                                index=y_train.index[24:])

    # 步骤2：对残差序列建立 VAR(p) 模型
    model = VAR(residuals_df, exog=exog_train[24:])
    results = model.fit(maxlags=lag_order, ic=None)

    # 步骤3：检查模型平稳性（特征值模）
    eigenvalues = np.linalg.eigvals(results.params.values[:6, :6])
    max_eigenvalue = np.max(np.abs(eigenvalues))
    print(f”最大特征值模: {max_eigenvalue:.4f}”)

    if max_eigenvalue >= 1.0:
        print(“警告：模型不满足平稳性条件！”)

    return results, residuals_df

def predict_hybrid_var(model, y_history, exog_test, baseline_lag=24):
    “””混合 VAR 预测：周期基线 + VAR 残差修正”””
    predictions = []
    residuals_history = []

    # 初始化残差历史（用于 VAR 预测）
    for i in range(len(y_history) - baseline_lag):
        res = y_history.iloc[i + baseline_lag] - y_history.iloc[i]
        residuals_history.append(res.values)

    for i in range(len(exog_test)):
        # 周期基线：前 24 小时同时刻值
        if i < baseline_lag:
            baseline = y_history.iloc[-(baseline_lag - i)]
        else:
            baseline = predictions[i - baseline_lag]

        # VAR 残差预测
        lag_order = model.k_ar
        recent_residuals = residuals_history[-lag_order:]
        residual_pred = model.forecast(y=np.array(recent_residuals),
                                       exog=exog_test[i:i+1],
                                       steps=1)[0]

        # 混合预测：基线 + 残差
        pred = baseline + residual_pred
        predictions.append(pred)
        residuals_history.append(residual_pred)

    return np.array(predictions)

def evaluate_model(y_true, y_pred, pollutant_names):
    “””评估模型性能：MAE、RMSE、R²”””
    from sklearn.metrics import mean_absolute_error, mean_squared_error, r2_score

    results = []
    for i, name in enumerate(pollutant_names):
        mae = mean_absolute_error(y_true[:, i], y_pred[:, i])
        rmse = np.sqrt(mean_squared_error(y_true[:, i], y_pred[:, i]))
        r2 = r2_score(y_true[:, i], y_pred[:, i])

        results.append({
            'pollutant': name,
            'MAE': mae,
            'RMSE': rmse,
            'R2': r2
        })

    return pd.DataFrame(results)
\end{lstlisting}

\subsubsection{稀疏 VAR 模型与系数分析}

\begin{lstlisting}[language=Python]
from sklearn.linear_model import MultiTaskLassoCV
import matplotlib.pyplot as plt
import seaborn as sns

def train_sparse_var(X_train, y_train, cv=5):
    “””稀疏 VAR：通过 L1 正则化筛选显著交叉滞后项”””
    # 使用交叉验证自动选择正则化参数
    model = MultiTaskLassoCV(cv=cv, max_iter=5000, random_state=42,
                             alphas=np.logspace(-4, 0, 50))
    model.fit(X_train, y_train)

    # 统计稀疏度
    coef_matrix = model.coef_
    strict_sparsity = np.sum(np.abs(coef_matrix) < 1e-10) / coef_matrix.size
    practical_sparsity = np.sum(np.abs(coef_matrix) < 0.01) / coef_matrix.size

    print(f”最优正则化参数 alpha: {model.alpha_:.6f}”)
    print(f”严格稀疏度（|coef| < 1e-10）: {strict_sparsity:.2%}”)
    print(f”实际稀疏度（|coef| < 0.01）: {practical_sparsity:.2%}”)

    return model

def extract_var_coefficient_matrices(model, feature_names, pollutant_names):
    “””提取 VAR 系数矩阵 A1 和 A2”””
    coef = model.coef_  # shape: (n_pollutants, n_features)

    # 找到滞后1和滞后2对应的特征索引
    lag1_indices = [i for i, name in enumerate(feature_names)
                    if any(p in name and 'lag1' in name for p in pollutant_names)]
    lag2_indices = [i for i, name in enumerate(feature_names)
                    if any(p in name and 'lag2' in name for p in pollutant_names)]

    # 提取系数矩阵
    A1 = coef[:, lag1_indices]  # shape: (6, 6)
    A2 = coef[:, lag2_indices]  # shape: (6, 6)

    return A1, A2

def visualize_coefficient_matrix(A, pollutant_names, title):
    “””可视化系数矩阵热图”””
    plt.figure(figsize=(8, 6))
    sns.heatmap(A, annot=True, fmt='.2f', cmap='RdBu_r',
                center=0, vmin=-5, vmax=5,
                xticklabels=pollutant_names,
                yticklabels=pollutant_names,
                cbar_kws={'label': '系数值'})
    plt.title(title, fontsize=14)
    plt.xlabel('来源污染物（t-k）', fontsize=12)
    plt.ylabel('目标污染物（t）', fontsize=12)
    plt.tight_layout()
    plt.savefig(f'{title}.pdf', dpi=300, bbox_inches='tight')
    plt.close()

def analyze_cross_lag_effects(A1, A2, pollutant_names, threshold=0.5):
    “””分析显著的交叉滞后效应”””
    print(“\n=== 显著交叉滞后效应分析 ===”)

    for lag, A in enumerate([A1, A2], start=1):
        print(f”\n滞后 {lag} 小时：”)
        for i, target in enumerate(pollutant_names):
            for j, source in enumerate(pollutant_names):
                if i != j and abs(A[i, j]) > threshold:
                    direction = “正向” if A[i, j] > 0 else “负向”
                    print(f”  {source} -> {target}: {A[i,j]:.3f} ({direction})”)
\end{lstlisting}

\subsubsection{协整检验与 VECM 实验}

\begin{lstlisting}[language=Python]
from statsmodels.tsa.vector_ar.vecm import coint_johansen, VECM

def johansen_cointegration_test(data, det_order=0, k_ar_diff=1):
    “””Johansen 协整检验：确定协整秩”””
    result = coint_johansen(data, det_order=det_order, k_ar_diff=k_ar_diff)

    # 提取特征值和迹统计量
    trace_stats = result.lr1  # 迹统计量
    critical_values = result.cvt  # 临界值（90%, 95%, 99%）

    print(“=== Johansen 协整检验结果 ===”)
    print(f”{'秩':<5} {'迹统计量':<12} {'5%临界值':<12} {'结论'}”)
    print(“-” * 50)

    # 确定协整秩
    coint_rank = 0
    for i in range(len(trace_stats)):
        conclusion = “拒绝” if trace_stats[i] > critical_values[i, 1] else “接受”
        print(f”{i:<5} {trace_stats[i]:<12.4f} {critical_values[i,1]:<12.4f} {conclusion}”)

        if trace_stats[i] > critical_values[i, 1]:
            coint_rank = i + 1

    print(f”\n协整秩: {coint_rank}”)

    return coint_rank, trace_stats, critical_values

def train_vecm_model(data, coint_rank, k_ar_diff=1):
    “””训练 VECM 模型”””
    if coint_rank == 0:
        print(“协整秩为0，无法建立 VECM 模型”)
        return None

    if coint_rank == data.shape[1]:
        print(“协整秩等于变量数（满秩），VECM 等价于 VAR”)

    model = VECM(data, k_ar_diff=k_ar_diff, coint_rank=coint_rank)
    results = model.fit()

    print(“\n=== VECM 模型摘要 ===”)
    print(f”协整秩: {coint_rank}”)
    print(f”滞后阶数: {k_ar_diff}”)
    print(f”AIC: {results.aic:.4f}”)
    print(f”BIC: {results.bic:.4f}”)

    return results

def compare_vecm_var(data, coint_rank):
    “””对比 VECM 与 VAR 性能”””
    from statsmodels.tsa.api import VAR

    # 训练 VECM
    vecm_model = train_vecm_model(data, coint_rank, k_ar_diff=2)

    # 训练 VAR(2)
    var_model = VAR(data).fit(maxlags=2, ic=None)

    print(“\n=== 模型对比 ===”)
    print(f”VECM AIC: {vecm_model.aic:.4f}”)
    print(f”VAR(2) AIC: {var_model.aic:.4f}”)
    print(f”VECM BIC: {vecm_model.bic:.4f}”)
    print(f”VAR(2) BIC: {var_model.bic:.4f}”)

    return vecm_model, var_model
\end{lstlisting}

\subsubsection{校园活动建议生成与前瞻性预警}

\begin{lstlisting}[language=Python]
def generate_daypart_recommendations(aqi_predictions, dates, temp_history):
    “””生成分时段活动建议（含前瞻性预警）”””
    recommendations = []

    for date in dates.unique():
        day_data = aqi_predictions[aqi_predictions['date'] == date]

        # 划分时段并取峰值 AQI
        morning_aqi = day_data[(day_data['hour'] >= 6) &
                               (day_data['hour'] <= 11)]['AQI'].max()
        afternoon_aqi = day_data[(day_data['hour'] >= 12) &
                                 (day_data['hour'] <= 17)]['AQI'].max()
        evening_aqi = day_data[(day_data['hour'] >= 18) &
                               (day_data['hour'] <= 23)]['AQI'].max()

        # 前瞻性预警：检测前12小时低温高湿条件
        warning = check_meteorological_warning(date, temp_history)

        # 生成建议
        recommendations.append({
            'date': date,
            'morning_aqi': morning_aqi,
            'afternoon_aqi': afternoon_aqi,
            'evening_aqi': evening_aqi,
            'morning_advice': get_activity_advice(morning_aqi, warning),
            'afternoon_advice': get_activity_advice(afternoon_aqi, warning),
            'evening_advice': get_activity_advice(evening_aqi, warning),
            'warning': warning
        })

    return pd.DataFrame(recommendations)

def check_meteorological_warning(date, temp_history, humidity_history):
    “””前瞻性预警：基于12小时滞后效应”””
    # 获取前12小时气象数据
    start_time = date - pd.Timedelta(hours=12)
    end_time = date
    temp_window = temp_history[start_time:end_time]
    humidity_window = humidity_history[start_time:end_time]

    # 判断条件：温度持续低于5°C且湿度高于80%
    low_temp = (temp_window < 5).sum() >= 8  # 至少8小时
    high_humidity = (humidity_window > 80).sum() >= 8

    if low_temp and high_humidity:
        return “前12小时低温高湿，预警午后颗粒物累积风险”
    else:
        return None

def get_activity_advice(aqi, warning=None):
    “””根据 AQI 等级返回活动建议”””
    # 基础建议
    if aqi <= 50:
        advice = “适宜户外活动；建议正常通风与锻炼。”
    elif aqi <= 100:
        advice = “适宜户外活动；敏感人群可减少极长时间运动。”
    elif aqi <= 150:
        advice = “减少长时间高强度户外活动，敏感人群建议室内。”
    elif aqi <= 200:
        advice = “建议室内锻炼，缩短户外停留时间。”
    elif aqi <= 300:
        advice = “避免户外锻炼，活动尽量安排在室内。”
    else:
        advice = “避免一切户外活动，室内采取防护措施。”

    # 叠加前瞻性预警
    if warning:
        advice += f” 【预警】{warning}，建议将户外活动延后至次日。”

    return advice

def export_recommendations_to_csv(recommendations, output_path):
    “””导出活动建议表”””
    recommendations.to_csv(output_path, index=False, encoding='utf-8-sig')
    print(f”活动建议已导出至: {output_path}”)
\end{lstlisting}

\subsubsection{滞后相关分析与可视化}

\begin{lstlisting}[language=Python]
def calculate_lag_correlation(pollutant, meteorology, max_lag=48):
    “””计算污染物与气象变量的滞后相关系数”””
    correlations = []

    for lag in range(max_lag + 1):
        # 滞后 k 小时的相关系数
        met_lagged = meteorology.shift(lag)
        valid_idx = ~(pollutant.isna() | met_lagged.isna())

        if valid_idx.sum() > 0:
            corr = pollutant[valid_idx].corr(met_lagged[valid_idx])
            correlations.append({'lag': lag, 'correlation': corr})

    return pd.DataFrame(correlations)

def find_optimal_lag(lag_corr_df):
    “””找到最强相关的滞后时间”””
    max_idx = lag_corr_df['correlation'].abs().idxmax()
    optimal_lag = lag_corr_df.loc[max_idx, 'lag']
    max_corr = lag_corr_df.loc[max_idx, 'correlation']

    return int(optimal_lag), max_corr

def plot_lag_correlation_heatmap(pollutants, meteorology, max_lag=48):
    “””绘制滞后相关热图”””
    pollutant_names = ['PM2.5', 'PM10', 'O3', 'NO2', 'SO2', 'CO']
    met_names = ['温度', '湿度']

    results = np.zeros((len(pollutant_names), len(met_names), max_lag + 1))

    for i, p_name in enumerate(pollutant_names):
        for j, m_name in enumerate(met_names):
            lag_corr = calculate_lag_correlation(
                pollutants[p_name], meteorology[m_name], max_lag)
            results[i, j, :] = lag_corr['correlation'].values

    # 绘制热图
    fig, axes = plt.subplots(2, 1, figsize=(12, 8))

    for j, m_name in enumerate(met_names):
        sns.heatmap(results[:, j, :], ax=axes[j],
                    cmap='RdBu_r', center=0, vmin=-0.6, vmax=0.6,
                    xticklabels=range(0, max_lag+1, 6),
                    yticklabels=pollutant_names,
                    cbar_kws={'label': '相关系数'})
        axes[j].set_title(f'污染物与{m_name}的滞后相关', fontsize=14)
        axes[j].set_xlabel('滞后时间（小时）', fontsize=12)
        axes[j].set_ylabel('污染物', fontsize=12)

    plt.tight_layout()
    plt.savefig('lag_correlation_heatmap.pdf', dpi=300, bbox_inches='tight')
    plt.close()
\end{lstlisting}

\subsubsection{消融实验与鲁棒性分析}

\begin{lstlisting}[language=Python]
def ablation_study_independent_vs_joint(X_train, y_train, X_test, y_test):
    “””消融实验：独立模型 vs 联合VAR模型”””
    from sklearn.linear_model import Ridge

    # 策略1：独立Ridge模型
    independent_predictions = []
    for i in range(y_train.shape[1]):
        model = Ridge(alpha=1.0)
        model.fit(X_train, y_train[:, i])
        pred = model.predict(X_test)
        independent_predictions.append(pred)

    independent_predictions = np.array(independent_predictions).T

    # 策略2：联合VAR模型（已训练）
    # joint_predictions = ... (从混合VAR获取)

    # 评估对比
    from sklearn.metrics import mean_absolute_error

    independent_mae = np.mean([
        mean_absolute_error(y_test[:, i], independent_predictions[:, i])
        for i in range(y_test.shape[1])
    ])

    print(f”独立Ridge模型平均MAE: {independent_mae:.4f}”)
    # print(f”联合VAR模型平均MAE: {joint_mae:.4f}”)
    # print(f”改进幅度: {(independent_mae - joint_mae) / independent_mae * 100:.2f}%”)

def robustness_test_meteorological_shock(df, model, threshold=5):
    “””鲁棒性分析：识别气象突变时段”””
    # 识别温度骤变时段
    temp_diff = df['温度'].diff(periods=6).abs()
    shock_indices = temp_diff[temp_diff > threshold].index

    print(f”识别到 {len(shock_indices)} 个气象突变时段”)

    # 在突变时段评估模型性能
    shock_mae = []
    for idx in shock_indices:
        if idx + 6 < len(df):
            window = df.loc[idx:idx+6]
            # 评估该窗口的预测误差
            # ...

    return shock_indices, shock_mae
\end{lstlisting}

\vspace{1em}
\noindent\textbf{说明：}以上代码片段展示了本文核心算法的实现逻辑。完整代码包含详细的参数设置、异常处理、可视化函数等，请参见支撑材料中的完整脚本文件。

\end{document}
